\documentclass[11pt]{article}



\usepackage[T1]{fontenc}
\usepackage[utf8]{inputenc}
\usepackage[english]{babel}
\usepackage[margin=1in]{geometry}
\usepackage{lmodern}
\usepackage{microtype}
\usepackage{mathtools}
\usepackage{amsmath,amssymb,amsthm}
\usepackage{array}
\usepackage{booktabs}
\usepackage{tabularx}
\usepackage{longtable}
\usepackage{float}
\usepackage{graphicx}
\usepackage{tikz}
\usetikzlibrary{arrows.meta,calc,fit,positioning}
\usepackage{pdflscape}
\usepackage{caption}
\usepackage{enumitem}
\usepackage{fancyhdr}
\usepackage[numbers]{natbib}
\usepackage{doi}
\usepackage{xurl}
\usepackage{hyperref}
\usepackage{cleveref}
\usepackage{orcidlink}
\usepackage{titling}

\newtheorem{theorem}{Theorem}[section]
\newtheorem{lemma}[theorem]{Lemma}
\newtheorem{proposition}[theorem]{Proposition}
\newtheorem{corollary}[theorem]{Corollary}
\theoremstyle{definition}
\newtheorem{definition}[theorem]{Definition}
\newtheorem{remark}[theorem]{Remark}
\newtheorem{warning}[theorem]{Warning}
\newtheorem{obligation}[theorem]{Obligation}
\newtheorem{nonclaim}[theorem]{Nonclaim}

\urlstyle{same}
\captionsetup{font=small, labelfont=bf, labelsep=period}
\captionsetup[table]{skip=6pt, position=top}
\captionsetup[figure]{skip=6pt}
\crefname{theorem}{theorem}{theorems}
\Crefname{theorem}{Theorem}{Theorems}
\crefname{lemma}{lemma}{lemmas}
\Crefname{lemma}{Lemma}{Lemmas}
\crefname{proposition}{proposition}{propositions}
\Crefname{proposition}{Proposition}{Propositions}
\crefname{corollary}{corollary}{corollaries}
\Crefname{corollary}{Corollary}{Corollaries}
\crefname{definition}{definition}{definitions}
\Crefname{definition}{Definition}{Definitions}
\crefname{remark}{remark}{remarks}
\Crefname{remark}{Remark}{Remarks}
\crefname{warning}{warning}{warnings}
\Crefname{warning}{Warning}{Warnings}
\crefname{obligation}{obligation}{obligations}
\Crefname{obligation}{Obligation}{Obligations}
\crefname{nonclaim}{nonclaim}{nonclaims}
\Crefname{nonclaim}{Nonclaim}{Nonclaims}
\setlist[enumerate]{leftmargin=*, itemsep=2pt, topsep=4pt}
\setlength{\droptitle}{-3.2em}
\pretitle{\begin{center}\LARGE\bfseries}
\posttitle{\par\end{center}\vspace{0.5em}}
\preauthor{\begin{center}\normalsize}
\postauthor{\par\end{center}\vspace{-0.4em}}
\predate{\begin{center}\small}
\postdate{\par\end{center}\vspace{0.6em}}
\clubpenalty=10000
\widowpenalty=10000
\displaywidowpenalty=10000
\interfootnotelinepenalty=10000
\allowdisplaybreaks[2]
\fancypagestyle{firstpage}{\fancyhf{}
  \fancyfoot[C]{\scriptsize Automorph Inc., Wilmington, DE, USA \quad \texttt{ioannis@automorph.io} \quad \textcopyright\ 2026 \quad CC-BY 4.0 \quad \doi{10.5281/zenodo.20713981}}
  \renewcommand{\headrulewidth}{0pt}
  \renewcommand{\footrulewidth}{0pt}
}



\newcommand{\Q}{\mathbb{Q}}
\newcommand{\R}{\mathbb{R}}
\newcommand{\Z}{\mathbb{Z}}
\newcommand{\Qp}{\mathbb{Q}_p}
\newcommand{\Zp}{\mathbb{Z}_p}
\newcommand{\Sha}{\operatorname{Sha}}
\newcommand{\Reg}{\operatorname{Reg}}
\newcommand{\Tam}{\operatorname{Tam}}
\newcommand{\tors}{\operatorname{tors}}
\newcommand{\tr}{\operatorname{tr}}
\newcommand{\Pf}{\operatorname{Pf}}
\newcommand{\Fitt}{\operatorname{Fitt}}
\newcommand{\cores}{\operatorname{cores}}
\newcommand{\rank}{\operatorname{rank}}
\newcommand{\Hf}{H^1_f}
\newcommand{\kapr}{\kappa_r}

\newcommand{\XiC}{\Xi_C(D\mid L)}        \newcommand{\KDD}{K_{DD}}
\newcommand{\KDL}{K_{DL}}
\newcommand{\KLL}{K_{LL}}
\newcommand{\Kdagger}{K_{LL}^{\dagger}}  \newcommand{\vsrc}{\mathrm{vsrc}}        \newcommand{\Jone}{J_1}
\newcommand{\Jtwo}{J_2}
\newcommand{\Jtwelve}{J_{12}}
\newcommand{\MCphi}{M_C^{\varphi}}

\setlength{\emergencystretch}{3.5em}
\sloppy

\title{Decomposing the Bloch--Kato Fundamental Line: Adequacy and No-Go
Results for the BSD Invariants}
\author{Ioannis Tsiokos\,\orcidlink{0009-0009-7659-5964}}
\date{31 August 2026}
\hypersetup{
  pdftitle={Decomposing the Bloch--Kato Fundamental Line: Adequacy and No-Go Results for the BSD Invariants},
  pdfauthor={Ioannis Tsiokos},
  hidelinks
}

\begin{document}
\maketitle
\thispagestyle{firstpage}

\begin{abstract}
For an elliptic curve \(E/\Q\), the Bloch--Kato fundamental line is often
read through a product of public scalar invariants.  This paper asks,
column by column, which typed source data such a scalar retains.  We
decompose the BSD/Bloch--Kato carrier into five typed columns:
height/regulator, analytic/period, finite/Tamagawa, \(p\)-adic, and
determinant-assembly.  Each column is equipped with a Schur-complement
adequacy residual \(\XiC\), which vanishes exactly when the admitted native
source acts as the identity on the readout; this is a finite,
mathlib-free linear-algebra criterion.  A comparison theorem contracts the
typed five-column atlas to the standard Bloch--Kato/Tamagawa scalar
package, with zero residual for that contraction, so the usual scalar is
recovered as a projection of a richer typed object and the
determinant-orientation coordinate lies in the kernel.  The accompanying
no-go audit suite records where public scalar rows underdetermine their
typed targets: finite scalar shadows, dimension-only \(\Sha\) data, the
Gram-matrix/regulator shadow, finite-window promotion, and finite-prime
coverage.  The results are finite typed statements; recognition sources
and classical arithmetic inputs such as Gross--Zagier, Iwasawa main
conjecture input, and Cassels--Tate comparison are carried as hypotheses,
not proved.  The paper does not prove BSD; its finite adequacy and no-go
claims are mechanized in a mathlib-free Lean development.
\end{abstract}

\medskip
\noindent\textbf{Keywords.} Birch--Swinnerton-Dyer conjecture; Bloch--Kato fundamental line; elliptic curves; Schur complement; regulator; Tate--Shafarevich group; formal verification (Lean); Six Birds Theory; emergence calculus.

\medskip
\noindent
Standing conventions: $E/\Q$ an elliptic curve, $M_E=h^1(E)(1)$,
$r=\rank E(\Q)$, $r_{\mathrm{an}}=\operatorname{ord}_{s=1}L(E,s)$. The Schur
adequacy residual is $\XiC=K_{DD}-K_{DL}K_{LL}^{\dagger}K_{LD}$; a column
readout is \emph{adequate} when $\XiC=0$. The five per-column source
transports are recognition obligations, supplied or imported and carried as
hypotheses rather than derived here.

\section{Introduction}\label{sec:introduction}


The Birch--Swinnerton-Dyer formula~\citep{BirchSwinnertonDyer1965,Tate1966bsd} is usually first met as a scalar
identity.  For an elliptic curve $E/\Q$ of rank $r$, the standard leading
coefficient form compares
\[
  \frac{L^{(r)}(E,1)}{r!}
  \quad\text{with}\quad
  \Omega_E\,\Reg(E/\Q)\,
  \frac{\#\Sha(E/\Q)\,\Tam(E)}{\#E(\Q)_{\tors}^{\,2}} .
\]
This product is already visibly composite: an analytic leading term and
period, a height regulator, a finite Tamagawa--torsion--$\Sha$ factor, local
$p$-adic data, and the determinant-line conventions that make all of these
terms live in a common orientation.  The question addressed in this paper is
not whether the scalar BSD identity is true.  It is the more local diagnostic
question: when one separates the Bloch--Kato fundamental line~\citep{BlochKato1990,FontainePerrinRiou1994} into these
typed pieces, which pieces carry enough native data to recover their intended
readout, and where is information lost by passing to a public scalar shadow?

The word ``typed'' is used here in its elementary sense.  We keep the five
classes of data in five labelled slots rather than immediately multiplying
them into one number.  The height/regulator slot contains the Mordell--Weil
lattice and the Neron--Tate pairing; the analytic/period slot contains the
leading Taylor coefficient and the real period; the finite slot contains
Tamagawa factors, torsion, the $p$-primary $\Sha$ data, and Cassels--Tate
pairings; the $p$-adic slot contains the ordinary or signed local data; and
the determinant slot contains the orientation, sign, and unit conventions
needed for a single determinant line.  This decomposition is bookkeeping
only in the innocuous sense that a direct-sum decomposition is: it is
precisely the bookkeeping needed to ask whether each column is adequate.

The adequacy test is a Schur-complement calculation~\citep{Zhang2005schur,HornJohnson2013}.  Suppose that $L$ is
the native family admitted in a column and $D$ is the readout one wants to
recover from that column.  Form the block Gram quantities
\[
  \KLL,\qquad \KDL,\qquad \KDD,
\]
where $\KLL$ records the covariance of the admitted source directions,
$\KDL$ records the pairing between the readout and those source directions,
and $\KDD$ records the readout's own pre-source energy.  The adequacy
residual is
\[
  \XiC=\KDD-\KDL\,\Kdagger\,K_{LD},
\]
with $\Kdagger$ the Moore--Penrose inverse of $\KLL$ in the finite
calculation at hand.  The column is called \emph{adequate} when
$\XiC=0$: after projecting the desired readout onto the admitted source
directions, no residual direction remains.  If $\XiC>0$ in a scalar trace
calculation, then the admitted row is not enough; it has forgotten a
direction that is present in the typed target.

All adequacy computations in this paper are finite and explicit.  The
linear algebra is deliberately elementary: finite index sets, row vectors,
dot products, identity matrices, matrix multiplication, and traces.  The
collapses are not appeals to a black-box functional-analytic statement.
They are the finite identity-matrix action
\[
  I v=v
\]
and its Schur-complement consequence.  For example, if the full native
height matrix is admitted, the readout row $d_H$ for
$d(\log\det)_H$ projects through the identity source, so
$d_Hd_H^{\mathsf T}-d_H I d_H^{\mathsf T}=0$.  If only a scalar determinant
is admitted, the orthogonal complement has positive trace.  The same
pattern recurs in the analytic/period, finite-source, ordinary $p$-adic,
signed $p$-adic, and rank-two height calculations.

This finite Schur-complement language is the specialization, for the BSD
fundamental line, of the adequacy calculus developed in the surrounding
Six Birds framework \citep{TsiokosNeedles2026,TsiokosFoundationsII2026}.
For the reader of the present paper, no framework machinery is needed
beyond the following translation.  An ``adequacy residual'' is simply the
Schur complement above.  A ``typed five-column decomposition'' is simply
the decision to keep the height, period, finite, $p$-adic, and determinant
data in separate labelled columns.  An ``audit-result no-go'' is a recorded
boundary statement: it identifies where the theorem-grade literature, in
the audited form used here, does not supply the missing identity.  The
interaction vocabulary and audit discipline are inherited from the finite
calculus of \citet{TsiokosFoundationsIII2026}; the mathematical statements
below are written directly in the language of the BSD factors.

A low-conductor example makes the separation concrete.  For the curve
$11a1$~\citep{Cremona1997,LMFDB}, the scalar rank-zero BSD comparison has
\[
  \frac{L(E,1)}{\Omega_E^+}=0.2
  \qquad\text{and}\qquad
  \frac{\prod_v c_v\,\#\Sha(E/\Q)}{\#E(\Q)_{\tors}^{\,2}}
  =\frac{5}{25}=0.2.
\]
As a scalar check this is a single equality.  In the five-column atlas it is
not a single source.  The analytic/period column contributes the quotient
$L(E,1)/\Omega_E^+$; the height/regulator column is rank-zero and therefore
vacuous; the finite column remembers that the number $5/25$ arose from a
Tamagawa factor and torsion rather than from Cassels--Tate pairing data;
the local $p$-adic columns ask which ordinary or signed source data are
present at each prime; and the determinant column records the orientation
needed to assemble the product.  The scalar equality is therefore compatible
with, but coarser than, the typed comparison.  The point of the apparatus is
to measure this coarsening column by column.

This is also the reason for working with readouts rather than only with
equalities of final products.  A scalar product can be correct for several
different reasons.  In the finite factor, a valuation can sit in the
Tamagawa part or in the $\Sha$ part and still contribute the same exponent
to the public scalar quotient.  In the signed supersingular local theory,
the two signed directions can have the same unsigned aggregate while
differing in their $+$ and $-$ components.  In rank at least two, a Gram
matrix can give the regulator once a Mordell--Weil basis is already known,
but the matrix itself does not name independent motivic-cycle sources.
These are not paradoxes.  They are exactly what one should expect when a
multi-column object is projected to a single row.  The Schur residual is the
linear-algebraic device that records which directions survive the projection
and which directions have been lost.

The paper therefore uses ``native family'' in a modest sense: it means the
data one is willing to admit as the source for a column before applying the
readout.  For the analytic/period column, the native family is the pair
$(A_E,\Omega_E^+)$, where
$A_E=L^{(r_{\mathrm{an}})}(E,1)/r_{\mathrm{an}}!$ is the analytic leading
coefficient and $\Omega_E^+$ is the real period, rather than the quotient
alone.  For the finite column, it is the tower containing Tamagawa, torsion,
$\Sha$, and Cassels--Tate
pairing data rather than the scalar finite factor alone.  For the signed
$p$-adic column, it is the signed pair of local directions rather than their
unsigned sum.  The adequacy theorem for a column says that this admitted
native family is large enough for the intended readout.  A no-go theorem
says that a smaller public row is not large enough, and it identifies an
explicit pair of inputs or an explicit residual direction witnessing the
loss.

Nothing in this setup requires the reader to adopt the surrounding
terminology as a black box.  The relevant matrices are small: identity
matrices of sizes $2$, $3$, $4$, and $5$, rank-one projectors, and diagonal
traces.  The Moore--Penrose notation is used only where the finite matrix
has an evident inverse or pseudoinverse in the displayed calculation.  The
rank-one residuals are ordinary rational computations such as
\[
  \tr(I_4-P_s)=4-1=3
\]
for the scalar finite row $s=(1,-2,1,0)$, or
\[
  \tr(I_3-P_{\mathrm{uns}})=3-2=1
\]
for the unsigned signed-local projection.  Thus the reader can check the
adequacy and inadequacy claims at the same level at which the statements are
made: finite-dimensional linear algebra over explicit coordinates.

The first part of the paper constructs the five columns.  The
Bloch--Kato component map sends the integrated BSD carrier to the
determinant-line target by the direct sum of the height/regulator,
analytic/period, finite, local $p$-adic, and determinant-assembly
components.  After choosing a comparison trivialization, the corresponding
defect decomposes additively into the five component defects.  Consequently,
if every component defect vanishes in the declared comparison, then the
total Bloch--Kato defect vanishes; contrapositively, a nonzero total defect
must be witnessed by at least one nonzero or legally incomparable component.

The order of the paper follows this diagnostic structure.  The
decomposition section fixes the five columns and the additive defect
equation.  The next four sections treat the height/regulator,
analytic/period, finite-source, and $p$-adic columns separately, always
asking the same question: does the declared native family collapse the
Schur residual for the declared readout?  The determinant-assembly section
then records the one remaining orientation coordinate needed to assemble
the local tensor into a global determinant-line object.  The comparison
section contracts the typed atlas back to the usual Bloch--Kato scalar
readout and proves that the typed source is adequate for that contraction.

The middle sections prove the adequacy collapses.  The height/regulator
column admits the full Neron--Tate pairing and collapses by the identity
action.  The analytic/period column admits both logarithmic coordinates
$(\log A_E,\log\Omega_E^+)$ and collapses for the readout
$\log A_E-\log\Omega_E^+$.  The finite-source column collapses only when
the Tamagawa, torsion, $\Sha$, and Cassels--Tate coordinates are all present;
the scalar finite factor and the dimension-only $\Sha$ row leave positive
residuals.  The ordinary $p$-adic row collapses with ordinary-control data
included, while omitting that coordinate leaves a one-dimensional residual.
The signed supersingular row collapses only after separating the $+$ and
$-$ directions; unsigned aggregation forgets their difference.  The
determinant-assembly calculation then shows a structural residual exactly
in the determinant-orientation coordinate when the four constructed
components are assembled without the final determinant row.

The comparison section contracts the typed five-column tensor to the usual
Bloch--Kato scalar package.  At the log-tangent level the comparison is the
row $(1,1,1,1,0)$: it reads the height, analytic, finite, and local
$p$-adic columns and kills the determinant-orientation coordinate.  Its
Schur residual is zero against the typed source, so the typed source is
adequate for the scalar Bloch--Kato readout.  At the same time, the kernel
contains the determinant coordinate, so the scalar package is a strict
one-dimensional shadow of the typed atlas.

Two calibrations complete the apparatus.  The first is the exterior algebra
$\vsrc$, generated in rank two by symbols $\Jone,\Jtwo$ and
$\Jtwelve=\Jone\Jtwo$ with $\Jone^2=\Jtwo^2=0$ and
$\Jtwo\Jone=-\Jtwelve$.  It is a determinant-shaped source algebra: the
wedge $\Jone\wedge\Jtwo$ records independence of two virtual motivic-cycle
directions without pretending that the corresponding motivic cycles have
been constructed.  The second is the normalization constant $\kapr$.  Once
the cascade convention $c_E=\Tam(E)$ is fixed, the cascade leading-term
form is equivalent to the standard Strong-BSD leading-coefficient form
exactly when
\[
  \kapr(E)=\frac{\#\Sha(E/\Q)}{\#E(\Q)_{\tors}^{\,2}}.
\]
This is an equivalence of forms, not a derivation of BSD.

The calibration sections have a different purpose from the column
sections.  The column sections test whether a given source row is adequate
for an already specified readout.  The $\vsrc$ section asks what sort of
formal source algebra would be capable of expressing rank-two independence
before any descent to actual motivic cycles is supplied.  The
$\kapr$ section asks how the cascade normalization must be chosen so that
the cascade leading-term expression is not a new conjecture but the
standard Strong-BSD expression written in the apparatus coordinates.  These
two calibrations constrain the apparatus at opposite ends: $\vsrc$ marks
the source-side independence problem, while $\kapr$ fixes the scalar
normalization at the target side.

The final mathematical contribution is the bounded no-go suite.  The
finite scalar factor is noninjective on the typed finite tower; the
dimension of $\Sha[p]$ is a shadow of Cassels--Tate determinant data; the
Neron--Tate Gram matrix does not by itself construct rank-$r$ motivic-cycle
sources; a five-curve window cannot promote to a global theorem depending
only on those five rows; and no fixed finite prime list certifies all local
support-prime rows.  These statements are useful because they are local and
typed: each identifies a precise missing direction rather than replacing
the missing direction by an undifferentiated failure of BSD.

The scope is deliberately narrow.  The no-go suite (NG-1--NG-5,
higher-rank, global, support-prime) records where theorem-grade Mode-A
literature (the established arithmetic results audited as external
provenance) does \textbf{not} reach; it is \textbf{not} a claim that the
missing identities cannot exist.  Adequacy collapses are proved at the
finite, mathlib-free, typed level (identity-matrix action; $\Q$ rank-1
projector traces); operator/motivic full generality is not claimed.  The
$\vsrc$ exterior-product relations are represented structurally on the
$2^r$ coordinate space; the dimension/calibration content is faithful,
full axiomatization is not claimed.  The five per-column transports are
recognition sources (carriers), not derived theorems.
 \section{Decomposition: the typed five columns}
\label{sec:decomposition}

The introduction separated the scalar BSD product into five kinds of data:
height/regulator, analytic/period, finite-source, local $p$-adic, and
determinant-assembly data.  We now make this separation into a typed object.
The point is not to replace the Bloch--Kato determinant line by five
unrelated objects, but to keep track of the five native sources before they
are assembled into the common determinant-line target.

Fix an elliptic curve $E/\Q$ and put $M_E=h^1(E)(1)$.  The five columns are
as follows.  The height/regulator column records the Mordell--Weil lattice
and the Neron--Tate determinant.  The analytic/period column records the
leading coefficient and period data.  The finite-source column records
Tamagawa factors, torsion, the $p$-primary $\Sha$ data, and Cassels--Tate
pairings.  For each prime $p$, the local $p$-adic column records the
ordinary or signed local source data.  The determinant-assembly column
records the orientation, sign, and unit conventions needed to compare all
four arithmetic columns in one determinant line.

The five-column map and the common Schur-residual test are summarized in
\Cref{tab:five-column-map,fig:five-column-schematic}.  The table fixes the
paper-prose names used in the later column sections; the figure records the
same bookkeeping as a source-to-target diagram.
\begin{table}[tbp]
\centering
\small
\caption{The five typed columns and their adequacy residuals.  Each row
separates the native source data from the determinant-line target before the
scalar BSD product is read out.}
\label{tab:five-column-map}
\begin{tabularx}{\textwidth}{@{}>{\raggedright\arraybackslash}p{0.18\textwidth}
>{\raggedright\arraybackslash}p{0.25\textwidth}
>{\raggedright\arraybackslash}p{0.22\textwidth}
>{\raggedright\arraybackslash}X@{}}
\toprule
Column & Carrier map & Residual & Adequate when \\
\midrule
Height/regulator &
\(\rho_{\mathrm{ht/reg}}^{\R}\) from the Mordell--Weil lattice and
Neron--Tate pairing &
\(\XiC\) for the trace of \(d(\log\det)_H\) &
the full height pairing is admitted, so the identity source kills the Schur
residual. \\
\addlinespace
Analytic/period &
\(\rho_{\mathrm{an/period}}^{\R}\) from \((A_E,\Omega_E^+)\) &
\(\XiC\) for \(\log A_E-\log\Omega_E^+\) &
both logarithmic coordinates are present; either coordinate alone leaves a
one-dimensional residual. \\
\addlinespace
Finite/Tamagawa &
\(\rho_{\mathrm{finite}}^{\R}\) from Tamagawa, torsion, \(\Sha[p^\infty]\),
and Cassels--Tate data &
\(\XiC\) for the typed finite-source readout &
the full finite-source tower, including the Cassels--Tate pairing, is
retained. \\
\addlinespace
\(p\)-adic &
\(\rho_p\) from ordinary or signed local data &
\(\XiC\) for the ordinary/signed local readout &
the appropriate ordinary or signed local source is admitted, not merely an
unsigned public shadow. \\
\addlinespace
Determinant assembly &
\(\rho_{\det}\) for orientation, sign, and unit compatibility &
\(\XiC\) for the assembled determinant comparison &
the determinant-orientation coordinate is included with the four arithmetic
columns. \\
\bottomrule
\end{tabularx}
\end{table}
 \begin{figure}[tbp]
\centering
\begin{tikzpicture}[
  font=\small,
  column/.style={draw, rounded corners=2pt, align=center, minimum width=2.6cm,
    minimum height=1.05cm, fill=black!3},
  residual/.style={draw, rounded corners=2pt, align=center, minimum width=7.5cm,
    minimum height=0.95cm, fill=white},
  target/.style={draw, rounded corners=2pt, align=center, minimum width=4.1cm,
    minimum height=1.05cm, fill=black!8},
  arrow/.style={-{Latex[length=2mm]}, thick}
]
\node[column] (ht) at (-5.6,1.7) {height\\regulator};
\node[column] (an) at (-2.8,1.7) {analytic\\period};
\node[column] (fin) at (0,1.7) {finite\\Tamagawa};
\node[column] (pad) at (2.8,1.7) {\(p\)-adic\\local};
\node[column] (det) at (5.6,1.7) {determinant\\assembly};
\node[residual] (xi) at (0,0) {Schur residual:\quad
  \(\XiC=K_{DD}-K_{DL}K_{LL}^{\dagger}K_{LD}\)};
\node[target] (bk) at (0,-1.55) {Bloch--Kato determinant-line target};

\foreach \n/\shift in {ht/-3cm,an/-1.5cm,fin/0cm,pad/1.5cm,det/3cm}
  \draw[arrow] (\n.south) .. controls +(0,-0.45) and +(0,0.45) ..
    ([xshift=\shift]xi.north);

\draw[arrow] (xi.south) -- (bk.north);
\node[align=center, font=\footnotesize] at (0,-2.6)
  {adequacy means the admitted native source leaves no Schur residual};
\end{tikzpicture}
\caption{The five native source columns are tested before they are assembled in
the common determinant-line target.}
\label{fig:five-column-schematic}
\end{figure}

The Bloch--Kato target is the determinant-line package attached to $M_E$~\citep{Deligne1987}:
the motive, its finite determinant line, the zeta element, and the chosen
trivialization.  The integrated BSD carrier is the corresponding typed
source object: it keeps the height/regulator, analytic/period,
finite-source, local $p$-adic, and determinant-assembly data in their
separate slots.  The first structural map of the paper is the
component-assembly map from this typed source to the determinant-line
target.

\begin{definition}[Bloch--Kato component-assembly map]
\label{def:apparatus:bk-component-map}
For $M_E=h^1(E)(1)$, let
\[
  C_{\mathrm{BSD}}(E)
  =
  \bigl(C_{\mathrm{ht}}(E),C_{\mathrm{an}}(E),C_{\mathrm{fin}}(E),
  \{C_p(E)\}_p,C_{\det}(E)\bigr)
\]
be the integrated BSD carrier with its typed height/regulator,
analytic/period, finite-source, local $p$-adic, and
determinant-assembly slots.  The source-side determinant slot
\(C_{\det}(E)\) is distinct from the Bloch--Kato target below.  Let
\[
  C_{\mathrm{BK}}(E)
  =
  \bigl(M_E,\Delta_f(M_E),z_{\mathrm{BK}}(M_E),\tau_{\mathrm{BK}}(M_E)\bigr)
\]
be the Bloch--Kato determinant-line target.  The component-assembly map is
the direct sum
\[
  \rho_{\mathrm{BK}}
  =
  \rho_{\mathrm{ht}}
  \oplus \rho_{\mathrm{an}}
  \oplus \rho_{\mathrm{fin}}
  \oplus \Bigl(\bigoplus_p \rho_p\Bigr)
  \oplus \rho_{\det}
  \;:\;
  C_{\mathrm{BSD}}(E)\longrightarrow C_{\mathrm{BK}}(E).
\]
Here $\rho_{\mathrm{ht}}$ sends the Mordell--Weil height and regulator
data to the global cohomological determinant side; $\rho_{\mathrm{an}}$
sends central leading-term and period data to the analytic
zeta/trivialization side; $\rho_{\mathrm{fin}}$ sends Tamagawa, torsion,
$\Sha[p^\infty]$, and Cassels--Tate data to the finite arithmetic
determinant side; $\rho_p$ sends the local $p$-adic analytic/Selmer data,
control data, and local normalization to the $p$-primary determinant side;
and $\rho_{\det}$ records the orientation, sign, and unit compatibility in
the same determinant line.
\end{definition}

The definition is a decomposition of the source side, not a theorem about
the truth of BSD.  Its role is to make the columns visible enough that the
adequacy residual of each column can be tested.  Once the columns are
visible, one may compare them through a common determinant-line
trivialization.  In additive coordinates, the comparison defect splits into
one contribution from each column.

\begin{definition}[Bridge-defect equation]
\label{def:apparatus:bridge-defect-equation}
After choosing a comparison trivialization of the determinant-line target,
write $\Delta_{\mathrm{BSD}}^{\mathrm{BK}}$ for the total Bloch--Kato defect
of the BSD bridge.  The bridge-defect equation is the additive
decomposition
\[
  \Delta_{\mathrm{BSD}}^{\mathrm{BK}}
  =
  E_{\mathrm{ht/reg}}
  +E_{\mathrm{an/period}}
  +E_{\mathrm{finite}}
  +\sum_p E_p
  +E_{\det}.
\]
The five summands are respectively the height/regulator defect, the
analytic/period defect, the finite-source defect, the local $p$-adic
defects, and the determinant-assembly defect.  Equivalently, in
multiplicative coordinates the same assertion says that the total component
ratio is the product of the corresponding five component ratios.
\end{definition}

\begin{theorem}[Component-killing]
\label{thm:apparatus:component-killing}
Assume that the five component defects are all legally comparable in the
chosen determinant-line trivialization.  If
\[
  E_{\mathrm{ht/reg}}=0,\qquad
  E_{\mathrm{an/period}}=0,\qquad
  E_{\mathrm{finite}}=0,\qquad
  E_p=0\ \text{for every }p,\qquad
  E_{\det}=0,
\]
then $\Delta_{\mathrm{BSD}}^{\mathrm{BK}}=0$.  Equivalently, if
$\Delta_{\mathrm{BSD}}^{\mathrm{BK}}\neq 0$, then at least one component
defect is nonzero or one of the asserted component comparisons is not
legally available.\footnote{Verified in Lean as
\texttt{SixBirdsBSD.Apparatus.Decomposition.componentKilling}; see App~\ref{app:formalization}.}
\end{theorem}

\begin{proof}
The comparison trivialization puts all five summands in the same additive
defect group.  Hence the bridge-defect equation of
Definition~\ref{def:apparatus:bridge-defect-equation} applies:
\[
  \Delta_{\mathrm{BSD}}^{\mathrm{BK}}
  =
  E_{\mathrm{ht/reg}}
  +E_{\mathrm{an/period}}
  +E_{\mathrm{finite}}
  +\sum_p E_p
  +E_{\det}.
\]
Under the stated hypotheses every summand on the right is zero.  The sum of
zero component defects is zero, so
$\Delta_{\mathrm{BSD}}^{\mathrm{BK}}=0$.

For the contraposition, suppose
$\Delta_{\mathrm{BSD}}^{\mathrm{BK}}\neq0$.  If all five component
comparisons were legally available and every component defect vanished, the
first part of the proof would force
$\Delta_{\mathrm{BSD}}^{\mathrm{BK}}=0$, a contradiction.  Therefore either
some declared component defect is nonzero, or the attempted componentwise
comparison is not legal in the determinant-line trivialization under
discussion.
\end{proof}

\begin{theorem}[Finite scalar factor is a public shadow]
\label{thm:apparatus:finite-scalar-public-shadow}
Let
\[
  C_{\mathrm{fin}}
  =
  C_{\mathrm{Tamagawa}}\oplus C_{\tors}\oplus
  C_{\Sha[p^\infty]}\oplus C_{\mathrm{CT}}
\]
be the typed finite-source tower.  The public scalar finite factor
\[
  \phi_{\mathrm{fin}}(C_{\mathrm{fin}})
  =
  \frac{|\Sha|\prod_\ell c_\ell}{|E(\Q)_{\tors}|^2}
\]
is not injective on this typed tower.  Equivalently, at a fixed prime $p$,
the valuation readout
\[
  q_p
  =
  v_p|\Sha|+v_p\Bigl(\prod_\ell c_\ell\Bigr)
  -2v_p|E(\Q)_{\tors}|
\]
does not determine the typed finite-source datum.  Thus a scalar
finite-factor row cannot be promoted to a finite-source carrier row without
additional source data.\footnote{Verified in Lean as
\texttt{SixBirdsBSD.Apparatus.Decomposition.finiteScalarPublicShadow}; see
App~\ref{app:formalization}.}
\end{theorem}

\begin{proof}
Fix a prime $p$ and consider only the three valuation coordinates appearing
in $q_p$: the $p$-primary part of $\Sha$, the Tamagawa product, and torsion.
The typed finite tower remembers which block carries a valuation, while the
scalar readout only remembers the integer combination
\[
  v_p|\Sha|+v_p\Bigl(\prod_\ell c_\ell\Bigr)
  -2v_p|E(\Q)_{\tors}|.
\]
Now compare two typed finite inputs.  In the first, put the valuation in the
$\Sha[p^\infty]$ block:
\[
  \bigl(v_p|\Sha|,\ v_p(\prod_\ell c_\ell),\
  v_p|E(\Q)_{\tors}|\bigr)=(1,0,0).
\]
In the second, put it in the Tamagawa block:
\[
  \bigl(v_p|\Sha|,\ v_p(\prod_\ell c_\ell),\
  v_p|E(\Q)_{\tors}|\bigr)=(0,1,0).
\]
The two typed inputs are distinct, since the first records a
$p$-primary $\Sha$ contribution and the second records a Tamagawa
contribution.  Nevertheless the scalar valuation is the same:
\[
  q_p=1+0-2\cdot0=1
  \qquad\text{and}\qquad
  q_p=0+1-2\cdot0=1.
\]
Thus the scalar map identifies two different typed finite-source rows.
Moreover the Cassels--Tate datum is absent from the displayed expression
for $q_p$ altogether.  Changing the Cassels--Tate pairing while holding the
three scalar valuation coordinates fixed leaves $q_p$ unchanged.  Hence the
scalar finite factor is a noninjective public shadow of the typed finite
tower.

This is an audit no-go in the sense used throughout the paper: it records
that the scalar finite row alone does not contain enough information to
recover the typed finite-source row.  It is not a claim that no stronger
arithmetic theorem, supplied with additional source data, can identify or
transport the missing finite information.
\end{proof}

\begin{nonclaim}[Decomposition scope]
\label{nonclaim:apparatus:decomposition-scope}
This section does not prove BSD, the Bloch--Kato conjecture, or ETNC~\citep{BurnsFlach2001}.  It
defines the five-column component map, proves the component-killing
consequence of the bridge-defect equation, and records the finite scalar
public-shadow no-go for the typed bridge.  The no-go is an audit result
about what the scalar row fails to determine; it is not a mathematical
non-existence claim about possible richer finite-source identities.  The
external transports from the five native columns to the determinant-line
target remain recognition sources carried separately, not derived theorems
of this decomposition section.
\end{nonclaim}
 \section{Height/regulator column}
\label{sec:height-regulator}

The height/regulator column is the part of the typed BSD carrier that keeps
the Mordell--Weil lattice and its Neron--Tate pairing~\citep{Neron1965,Silverman2009} before any comparison
with the global determinant line is made.  In classical BSD notation this is
the source of the regulator.  In the five-column language of
\Cref{sec:decomposition}, the native object is the real height pairing on
the free Mordell--Weil lattice, while the readout is the corresponding
cohomological determinant factor.

\begin{definition}[Real height-regulator map]\label{def:apparatus:height-regulator-map}
Let
\[
  \Lambda_E=E(\Q)/E(\Q)_{\tors},
  \qquad
  V_E=\Lambda_E\otimes_{\Z}\R,
\]
and let $\langle\cdot,\cdot\rangle_{\mathrm{NT}}$ be the Neron--Tate
pairing on $V_E$.  Write
\[
  h_{\mathrm{NT}}:V_E\longrightarrow V_E^*,
  \qquad
  P\longmapsto \langle P,\cdot\rangle_{\mathrm{NT}},
\]
for the associated linear map to the dual space.  The real
height-regulator map is
\[
  \rho_{\mathrm{ht/reg}}^{\R}:(V_E,\Lambda_E,\langle\cdot,\cdot\rangle_{\mathrm{NT}})
  \longmapsto \det(h_{\mathrm{NT}})\in\det(V_E^*)\otimes\det(V_E)^*,
\]
where the determinant is taken as an element of the indicated determinant
line.  If $(P_1,\ldots,P_r)$ is a $\Z$-basis of $\Lambda_E$ and
\[
  H_E=\bigl(\langle P_i,P_j\rangle_{\mathrm{NT}}\bigr)_{1\leq i,j\leq r},
\]
then the scalar coefficient of this determinant element in the induced
basis is
\[
  \det H_E=\Reg(E/\Q),
\]
with the empty determinant equal to $1$ when $r=0$.  This coefficient is
basis-independent: if the basis is changed by
$A\in\mathrm{GL}_r(\Z)$, then $H_E$ is replaced by
$A^{\mathsf T}H_EA$, and
\[
  \det(A^{\mathsf T}H_EA)
  =
  \det(A)^2\det(H_E)
  =
  \det(H_E),
\]
since $\det(A)=\pm1$.
\end{definition}

\begin{theorem}[Height tangent Schur collapse]\label{thm:apparatus:height-schur-collapse}
On the tangent height carrier $S_r=\operatorname{Sym}_r(\R)$, write
$x=\operatorname{vech}(\delta H)$ for the coordinate vector of a symmetric
height variation.  Take the audit energy to be the identity matrix $C=I$.
With the full native source $L_{\mathrm{full}}(x)=x$ and with readout
\[
  D_H(\delta H)
  =
  d(\log\det)_H(\delta H)
  =
  \tr(H^{-1}\delta H)
  =
  d_Hx,
\]
one has
\[
  \XiC=d_Hd_H^{\mathsf T}-d_H I d_H^{\mathsf T}=0 .
\]
The pre-source residual is $K_{DD}=\|d_H\|^2$ (e.g.\ $37a1$: $382.79\ldots$);
rank-$0$ rows are vacuous ($K_{DD}=0$).\footnote{Verified in Lean as
\texttt{SixBirdsBSD.Apparatus.HeightRegulator.heightSchurCollapse}; see App~\ref{app:formalization}.}
\end{theorem}

\begin{proof}
Let $n=\dim S_r=r(r+1)/2$ and identify a tangent vector with its
half-vectorized coordinate vector $x\in\R^n$.  The differential of
$\log\det$ at $H$ is the linear functional
\[
  x\longmapsto d_Hx,
\]
where $d_H$ is a row vector with one coordinate for each entry of
$\operatorname{vech}(\delta H)$.  Before admitting the height source, the
readout variance is simply
\[
  K_{DD}=d_Hd_H^{\mathsf T}
  =
  \sum_{i=1}^n (d_H)_i^2
  =
  \|d_H\|^2.
\]

Now admit the full tangent height carrier as the source.  Since
$L_{\mathrm{full}}(x)=x$ and $C=I$, the source block is
$K_{LL}=I$, its Moore--Penrose inverse is again $I$, and the cross blocks
are $K_{DL}=d_H$ and $K_{LD}=d_H^{\mathsf T}$.  The only computation is the
identity-matrix action.  If $v=(v_1,\ldots,v_n)^{\mathsf T}$, then the
$i$-th coordinate of $Iv$ is
\[
  \sum_{j=1}^n \delta_{ij}v_j=v_i,
\]
so $Iv=v$.  Applying this to $v=d_H^{\mathsf T}$ gives
$I d_H^{\mathsf T}=d_H^{\mathsf T}$, and therefore
\[
  K_{DL}K_{LL}^{\dagger}K_{LD}
  =
  d_H I d_H^{\mathsf T}
  =
  d_Hd_H^{\mathsf T}.
\]
Substituting into the Schur-complement residual gives
\[
  \XiC
  =
  K_{DD}-K_{DL}K_{LL}^{\dagger}K_{LD}
  =
  d_Hd_H^{\mathsf T}-d_Hd_H^{\mathsf T}
  =
  0.
\]

If $r=0$, then $S_0$ is the zero tangent space, $n=0$, and the row vector
$d_H$ has no coordinates.  The empty sum defining
$d_Hd_H^{\mathsf T}$ is $0$, so $K_{DD}=0$ already before any source is
admitted.  Thus rank-$0$ rows are vacuous in exactly the stated sense.
\end{proof}

\begin{obligation}[Refined height-regulator BK transport]\label{obl:apparatus:rho-ht-reg-transport}
The full Bloch--Kato height/regulator component also requires a comparison
transport
\[
  \beta_{\mathrm{BK,ht}}:
  \det(V_E^*)\otimes\det(V_E)^*
  \longrightarrow
  \Delta_f(M_E)_{\mathrm{ht/reg}}.
\]
This transport identifies the Neron--Tate determinant with the
Beilinson--Deligne regulator normalization~\citep{Beilinson1985,Schneider1982}, including the coefficient
convention, the archimedean Deligne-cohomology normalization, the determinant
orientation and trivialization, and compatibility with the finite and
determinant-assembly columns without absorbing them.  In this paper it is a
recognition-source carrier: the comparison is named as external arithmetic
content and carried separately, not derived from the finite typed
height-column calculation.
\end{obligation}

\begin{warning}[Carrier discipline for $571a1$]\label{warn:apparatus:571a1-carrier}
The curve $571a1$ has Mordell--Weil rank $0$ with nontrivial
$\Sha[2]$/Cassels--Tate data.  Its lawful height/regulator row is therefore
the rank-$0$ row described above: no nonzero height tangent is present in
this column.  Rank-$2$ height/regulator examples in this paper use $389a1$.
\end{warning}

 \section{Analytic/period column}
\label{sec:analytic-period}

The analytic/period column keeps the analytic leading coefficient and the
real period before they are mixed with height, finite, or local data.  Thus
the native pair is $(A_E,\Omega_E^+)$, where
$A_E=L^{(r_{\mathrm{an}})}(E,1)/r_{\mathrm{an}}!$ is the standard BSD
leading coefficient and $\Omega_E^+$ is the real period.  In log-tangent
coordinates the readout is the quotient coordinate
$\log A_E-\log\Omega_E^+$.

\begin{definition}[Real analytic/period map]\label{def:apparatus:analytic-period-map}
By modularity~\citep{Wiles1995,BreuilConradDiamondTaylor2001} $L(E,s)=L(f_E,s)$.  Let
\[
  A_E=\frac{L^{(r_{\mathrm{an}})}(E,1)}{r_{\mathrm{an}}!}
\]
be the analytic leading coefficient, equivalently the leading Taylor
coefficient of $L(E,s)$ at $s=1$, and let
\[
  \Omega_E^+=\int_{\gamma_\infty}\omega_E
\]
be the real period.  The native analytic/period family is
\[
  L_{\mathrm{an/per}}(E)=(A_E,\Omega_E^+).
\]
Its isolated scalar readout is $A_E/\Omega_E^+$, and the real
analytic/period map is
\[
  \rho_{\mathrm{an/period}}^{\R}:(A_E,\Omega_E^+)\longmapsto
  (A_E/\Omega_E^+)\,e_{\mathrm{an/per}}
  \in L_{\mathrm{an/per}}(E)_{\R}.
\]
This native column contains only analytic leading-term and period data:
height, torsion, Tamagawa, and $\Sha$ data are excluded.
\end{definition}

\begin{theorem}[Analytic/period tangent Schur collapse]\label{thm:apparatus:analytic-period-schur-collapse}
In the log-tangent chart
\[
  x=(x_A,x_\Omega)=(\log A_E,\log\Omega_E^+)\in\R^2,
\]
take $C=I_2$, take the full native source to be
$L_{\mathrm{full}}(x)=x$, and take the dissolving readout to be
\[
  D(x)=x_A-x_\Omega.
\]
Equivalently, $D(x)=dx$ for the row vector $d=(1,-1)$.  Then
$\KLL=I_2$, $\KDD=dd^{\mathsf T}=2$, and
\[
  \XiC=2-dI_2d^{\mathsf T}=0.
\]
With no source admitted the residual is $2$; with only the $x_A$ coordinate
admitted the residual is $1$; with both native coordinates admitted the
residual is $0$.\footnote{Verified in Lean as
\texttt{SixBirdsBSD.Apparatus.AnalyticPeriod.analyticPeriodSchurCollapse}; see App~\ref{app:formalization}.}
\end{theorem}

\begin{proof}
The readout row is $d=(1,-1)$, so before admitting any source its readout
energy is
\[
  \KDD=dd^{\mathsf T}=1^2+(-1)^2=2.
\]
Now admit the full analytic/period source
$L_{\mathrm{full}}(x)=x$ with audit energy $C=I_2$.  Then the source block
is $\KLL=I_2$, its Moore--Penrose inverse is again
$\Kdagger=I_2$, and the cross blocks are $\KDL=d$ and
$K_{LD}=d^{\mathsf T}$.  The identity matrix acts on a column vector
$v=(v_A,v_\Omega)^{\mathsf T}$ by
\[
  I_2v
  =
  \begin{pmatrix}1&0\\0&1\end{pmatrix}
  \begin{pmatrix}v_A\\v_\Omega\end{pmatrix}
  =
  \begin{pmatrix}v_A\\v_\Omega\end{pmatrix}
  =
  v.
\]
Applying this to $v=d^{\mathsf T}$ gives
$I_2d^{\mathsf T}=d^{\mathsf T}$.  Hence
\[
  \KDL\Kdagger K_{LD}
  =
  dI_2d^{\mathsf T}
  =
  dd^{\mathsf T}
  =
  2.
\]
The Schur-complement residual is therefore
\[
  \XiC
  =
  \KDD-\KDL\Kdagger K_{LD}
  =
  2-2
  =
  0.
\]

The two orientation numbers in the statement are the same finite
calculation with smaller source.  With no source, no projection is removed,
so the residual is $\KDD=2$.  If only $x_A$ is admitted, the row
$d=(1,-1)$ decomposes orthogonally as the admitted part $(1,0)$ plus the
unseen part $(0,-1)$; the unseen part has squared length $1$.  Thus the
one-coordinate source leaves residual $1$, while the two-coordinate source
leaves none.
\end{proof}

\begin{obligation}[Refined analytic/period BK transport]\label{obl:apparatus:rho-an-period-transport}
The full Bloch--Kato analytic/period component requires a comparison
transport
\[
  \beta_{\mathrm{BK,an}}:
  L_{\mathrm{an/per}}(E)_{\R}
  \longrightarrow
  \Delta_f(M_E)_{\mathrm{an/per}}.
\]
This transport identifies
\[
  \frac{L^{(r_{\mathrm{an}})}(E,1)}
       {r_{\mathrm{an}}!\,\Omega_E^+}
\]
with the analytic zeta-element factor, including the rational-versus-integral
normalization, the Betti--de Rham comparison and real orientation, the
positive-rank leading-term convention, and separation from the other four
columns.  The rank-$0$ comparison is classical; in positive rank the
comparison is named external arithmetic content.  In this paper it is a
recognition-source carrier, not a theorem derived from the finite
analytic/period Schur calculation.
\end{obligation}

 \section{Finite-source column and the finite no-gos}
\label{sec:finite-source}

The finite-source column keeps the arithmetic finite data before it is
contracted to the public scalar finite factor.  Its typed tower separates
Tamagawa factors~\citep{Silverman1994}, torsion, the $p$-primary \(\Sha\) data, and the
Cassels--Tate pairing.  The point of this column is that the determinant
readout is pairing-aware: the Cassels--Tate form~\citep{PoonenStoll1999,Cassels1962} is source data, not a
number to be recovered from the scalar finite product after the fact.

\begin{definition}[Real finite-source map]\label{def:apparatus:finite-source-map}
The finite-source tower is
\[
  C_{\mathrm{fin}}(E)
  =
  C_{\mathrm{Tam}}\oplus C_{\tors}\oplus
  C_{\Sha[p^\infty]}\oplus C_{\mathrm{CT}}.
\]
Its native family is
\[
  L_{\mathrm{finite}}(E)
  =
  \bigl((c_v(E))_{v\mid N},\,E(\Q)_{\tors},\,
  (\Sha(E/\Q)[p^\infty])_p,\,\langle\cdot,\cdot\rangle_{\mathrm{CT}}\bigr).
\]
The real finite-source map is the tensor
\[
  \rho_{\mathrm{finite}}^{\R}
  =
  \rho_{\mathrm{Tam}}^{\R}
  \otimes
  \rho_{\tors}^{\R}
  \otimes
  \bigotimes_p \rho_{\Sha\text{-}\mathrm{CT},p}^{\R}.
\]
The Tamagawa factor is
\[
  \rho_{\mathrm{Tam}}^{\R}:
  (c_v)_{v\mid N}
  \longmapsto
  \bigotimes_{v\mid N} c_v e_v^{\mathrm{loc}},
  \qquad
  c_v=|\Phi_v(E)|,\quad
  \Phi_v=E(\Q_v)/E^0(\Q_v).
\]
The torsion factor is
\[
  \rho_{\tors}^{\R}:
  T(E)
  \longmapsto
  \bigotimes_p |T_p(E)|^{-2}e_p^{\tors}.
\]
For the \(\Sha\)--Cassels--Tate factor, put
\[
  S_p=\Sha[p^\infty]/\Sha_{\mathrm{div}}[p^\infty],
\]
equipped with its perfect alternating Cassels--Tate pairing
\(\langle\cdot,\cdot\rangle_{\mathrm{CT},p}\).  Then
\[
  \rho_{\Sha\text{-}\mathrm{CT},p}^{\R}:
  (S_p,\langle\cdot,\cdot\rangle_{\mathrm{CT},p})
  \longmapsto
  \Pf(\langle\cdot,\cdot\rangle_{\mathrm{CT},p})
  e_p^{\Sha\text{-}\mathrm{CT}}.
\]
For a trivial \(S_p\), the Pfaffian contribution is the unit.  The source
datum in this block is the pair
\((S_p,\langle\cdot,\cdot\rangle_{\mathrm{CT},p})\), not merely
\(|S_p|\) and not merely \(\dim_{\mathbb F_p}S_p[p]\).
\end{definition}

\begin{theorem}[Pairing-aware finite Schur collapse]\label{thm:apparatus:finite-source-schur-collapse}
At a support prime $p$, use finite tangent coordinates
\[
  x=(x_T,x_{\tors},x_{\Sha},x_{\mathrm{CT}})\in\R^4,
\]
where \(x_T=v_p(\prod_v c_v)\), \(x_{\tors}=v_p|E(\Q)_{\tors}|\),
\(x_{\Sha}=v_p|S_p|\), and \(x_{\mathrm{CT}}\) is the
Cassels--Tate-pairing coordinate.  With the pairing-aware full readout
\[
  D_{\mathrm{full}}(x)=x,
  \qquad
  L_{\mathrm{full}}(x)=x,
  \qquad
  C=I_4,
\]
the Schur residual collapses:
\[
  \XiC
  =
  I_4-I_4I_4^{\dagger}I_4
  =
  0.
\]
Thus the full typed finite tower is adequate for the pairing-aware finite
readout.\footnote{Verified in Lean as
\texttt{SixBirdsBSD.Apparatus.FiniteSource.finiteSourceSchurCollapse}; see App~\ref{app:formalization}.}
\end{theorem}

\begin{proof}
Because the full readout and the full source are both the identity on
\(\R^4\), the readout block is
\[
  \KDD=I_4.
\]
The source block is also \(\KLL=I_4\).  Since the Moore--Penrose inverse of
the identity matrix is the identity matrix, \(\Kdagger=I_4\).  The two cross
blocks are \(\KDL=I_4\) and \(K_{LD}=I_4\).

It remains to compute the projection term rather than name it.  For any
vector \(v=(v_1,v_2,v_3,v_4)^{\mathsf T}\),
\[
  I_4v=v.
\]
Equivalently, in coordinates,
\[
  (I_4v)_i=\sum_{j=1}^4\delta_{ij}v_j=v_i.
\]
Applying the same identity-matrix action column by column to \(I_4\) gives
\[
  I_4I_4=I_4.
\]
Therefore
\[
  \KDL\Kdagger K_{LD}
  =
  I_4I_4I_4
  =
  I_4.
\]
Substitution into the Schur-complement residual gives
\[
  \XiC
  =
  \KDD-\KDL\Kdagger K_{LD}
  =
  I_4-I_4
  =
  0.
\]
This calculation uses all four finite coordinates.  The scalar and
dimension-only shadows treated below are different readouts: they forget
part of the typed finite tower and therefore leave positive residuals.
\end{proof}

\begin{theorem}[Scalar finite projection (NG-1, Schur form)]\label{thm:apparatus:scalar-finite-shadow}
Let \(x=(x_T,x_{\tors},x_{\Sha},x_{\mathrm{CT}})\in\R^4\), and let the
scalar finite projection be
\[
  q=s\cdot x,
  \qquad
  s=(1,-2,1,0).
\]
The orthogonal rank-one projector onto the scalar row is
\[
  P_s=s^{\mathsf T}(ss^{\mathsf T})^{-1}s,
\]
and the residual against the full typed finite readout is
\[
  \Xi_{\mathrm{scalar}}=I_4-P_s,
  \qquad
  \tr(\Xi_{\mathrm{scalar}})=3>0.
\]
Moreover the scalar map is noninjective on the typed finite tower: the typed
finite inputs \((0,0,2,H)\) and \((2,0,0,0)\) both give \(q=2\), but place
the valuation in the \(\Sha\)--Cassels--Tate block and the Tamagawa block,
respectively.\footnote{Verified in Lean as
\texttt{SixBirdsBSD.Apparatus.FiniteSource.scalarFiniteShadow}; see App~\ref{app:formalization}.}
\end{theorem}

\begin{proof}
First compute the Schur residual.  The squared length of the scalar row is
\[
  ss^{\mathsf T}
  =
  1^2+(-2)^2+1^2+0^2
  =
  6.
\]
Thus the rank-one projector has entries
\[
  (P_s)_{ij}=\frac{s_i s_j}{6}.
\]
Its trace is the sum of its diagonal entries:
\[
  \tr(P_s)
  =
  \sum_{i=1}^4 \frac{s_i^2}{6}
  =
  \frac{1^2+(-2)^2+1^2+0^2}{6}
  =
  \frac{6}{6}
  =
  1.
\]
Since \(\tr(I_4)=4\), the residual trace is
\[
  \tr(\Xi_{\mathrm{scalar}})
  =
  \tr(I_4-P_s)
  =
  \tr(I_4)-\tr(P_s)
  =
  4-1
  =
  3.
\]
This is positive, so the scalar row leaves a three-dimensional residual
relative to the full four-coordinate finite source.

The same failure is visible without linear algebra.  Evaluate \(q\) on the
two displayed typed inputs.  For the \(\Sha\)--Cassels--Tate input
\((0,0,2,H)\),
\[
  q=1\cdot0-2\cdot0+1\cdot2+0\cdot H=2.
\]
For the Tamagawa input \((2,0,0,0)\),
\[
  q=1\cdot2-2\cdot0+1\cdot0+0\cdot0=2.
\]
The scalar value is the same, but the typed data are different: one input
records the contribution in the \(\Sha\)--Cassels--Tate block, while the
other records it in the Tamagawa block.  The scalar projection also ignores
the Cassels--Tate coordinate \(H\) itself, since the fourth coefficient of
\(s\) is zero.  Hence the scalar finite row is a public shadow of the typed
finite tower.  This is an audit no-go: it says the scalar row alone does not
recover the typed source; it is not a claim that no richer theorem can
supply the missing finite data.
\end{proof}

\begin{theorem}[{$\dim\Sha[p]$ projection (NG-2, Schur form)}]\label{thm:apparatus:dim-sha-shadow}
On the \(\Sha\)--Cassels--Tate subcarrier \(y=(m,\theta)\in\R^2\), where
\(m\) records the \(p\)-power elementary-divisor length and \(\theta\)
records the Cassels--Tate coordinate, the dimension-only projection
\[
  L_{\dim}(y)=(1,0)y
\]
leaves
\[
  \Xi_{\dim}
  =
  I_2-(1,0)^{\mathsf T}(1,0)
  =
  \operatorname{diag}(0,1),
  \qquad
  \tr(\Xi_{\dim})=1>0.
\]
It is noninjective on Cassels--Tate data: \((\Z/2\Z)^2\) with
\(\langle e,f\rangle=\tfrac12\) and \((\Z/4\Z)^2\) with
\(\langle e,f\rangle=\tfrac14\) have the same
\(\mathbb F_2\)-dimension of their \(2\)-torsion but different
Cassels--Tate determinant factors.\footnote{Verified in Lean as
\texttt{SixBirdsBSD.Apparatus.FiniteSource.dimShaShadow}; see App~\ref{app:formalization}.}
\end{theorem}

\begin{proof}
Let \(e=(1,0)\).  The dimension-only source keeps the \(m\)-coordinate and
forgets the \(\theta\)-coordinate.  Its projector is
\[
  e^{\mathsf T}e
  =
  \begin{pmatrix}1&0\\0&0\end{pmatrix}.
\]
Therefore the residual projector is
\[
  \Xi_{\dim}
  =
  I_2-e^{\mathsf T}e
  =
  \begin{pmatrix}1&0\\0&1\end{pmatrix}
  -
  \begin{pmatrix}1&0\\0&0\end{pmatrix}
  =
  \begin{pmatrix}0&0\\0&1\end{pmatrix}
  =
  \operatorname{diag}(0,1).
\]
Taking traces gives
\[
  \tr(\Xi_{\dim})=0+1=1>0.
\]
Thus the dimension-only row leaves exactly the Cassels--Tate direction
unexplained.

The witness shows that this residual is not a bookkeeping artifact.  The
group \(U_1=(\Z/2\Z)^2\) has \(U_1[2]=U_1\), so
\(\dim_{\mathbb F_2}U_1[2]=2\).  The group \(U_2=(\Z/4\Z)^2\) has
\[
  U_2[2]\cong(\Z/2\Z)^2,
\]
so \(\dim_{\mathbb F_2}U_2[2]=2\) as well.  The dimension readout therefore
cannot distinguish \(U_1\) from \(U_2\).

Their Cassels--Tate determinant factors differ.  For a two-generator
alternating pairing with matrix
\[
  \begin{pmatrix}0&a\\-a&0\end{pmatrix},
\]
the Pfaffian is \(a\).  The pairings in the statement have \(a=\tfrac12\)
and \(a=\tfrac14\), respectively, so they give different Pfaffian
contributions to the determinant-line trivialization.  Hence the dimension
\(\dim_{\mathbb F_2}\Sha[2]\) is a noninjective public shadow of the
\(\Sha\)--Cassels--Tate source.  As above, this is an audit no-go about a
specific public readout, not a non-existence claim about all possible
arithmetic refinements.
\end{proof}

\begin{obligation}[Refined finite-source BK transport]\label{obl:apparatus:rho-finite-transport}
The full Bloch--Kato finite component requires a comparison transport
\[
  \beta_{\mathrm{BK,finite}}:
  L_{\mathrm{finite}}(E)_{\R}
  \longrightarrow
  \Delta_f(M_E)_{\mathrm{finite}}.
\]
This transport fixes the integral local-condition conventions at bad primes,
the squared-torsion normalization, the Cassels--Tate Pfaffian and orientation
conventions, and compatibility with the determinant-assembly column.  In
this paper it is a recognition-source carrier: the finite-source transport
is named external arithmetic content and carried separately, not derived
from the finite typed Schur calculation.
\end{obligation}

 \section{p-adic column}
\label{sec:p-adic}

The local \(p\)-adic column records the local source data before it is
contracted to a public local scalar.  At an ordinary prime this means keeping
the ordinary-control coordinate together with the finite local condition and
dual exponential.  At a supersingular prime it means keeping the signed
\(+\) and \(-\) local directions separately, rather than replacing them by
their unsigned aggregate.

\begin{definition}[Real p-adic map]\label{def:apparatus:p-adic-map}
Let \(T=T_pE\) and \(V=V_p(E)=T\otimes\Qp\), viewed as a representation of
\(G_{\Qp}\).  The finite local condition~\citep{BlochKato1990,Fontaine1994} is
\[
  \Hf(\Qp,V)
  =
  \ker\bigl(H^1(\Qp,V)\to H^1(\Qp,V\otimes B_{\mathrm{cris}})\bigr),
\]
and the dual exponential is
\[
  \exp^*:\Hf(\Qp,V)\longrightarrow
  D_{\mathrm{dR}}(V)/\operatorname{Fil}^0D_{\mathrm{dR}}(V).
\]

If \(E\) has good ordinary reduction at \(p\)~\citep{Mazur1972}, so \(a_p\) is a \(p\)-unit
and the ordinary filtration is
\[
  0\longrightarrow T^+\longrightarrow T\longrightarrow T^-\longrightarrow0,
\]
then the native ordinary family is
\[
  L_p^{\mathrm{ord}}(E)=(T^+,\Hf(\Qp,V),\exp^*),
\]
and the real ordinary local map is
\[
  \rho_{p,\mathrm{ord}}^{\R}:
  (T^+,\Hf(\Qp,V),\exp^*)
  \longmapsto
  e_{T^+}\otimes\det(\exp^*)\,e_{\mathrm{loc},p}
  \in \Delta_f(V,T)_{p,\R}^{\mathrm{ord}}.
\]
If \(E\) has good supersingular reduction at \(p\)~\citep{Kobayashi2003,Pollack2003}, with \(a_p\equiv0\),
then the signed native family is
\[
  L_p^{\pm}(E)=(H^1_+,H^1_-,\exp^*),
\]
and the signed real local map is
\[
  \rho_{p,\pm}^{\R}:
  (H^1_+,H^1_-,\exp^*)
  \longmapsto
  e_+\otimes e_-\otimes\det(\exp^*)\,e_{\mathrm{loc},p}^{\pm}.
\]
\end{definition}

\begin{theorem}[Ordinary local Schur collapse]\label{thm:apparatus:ordinary-schur-collapse}
For a good-ordinary row, let
\[
  x=(x_{\mathrm{ord}},x_f,x_{\exp})\in\R^3
\]
be the ordinary-control, finite-local-condition, and dual-exponential
tangent coordinates.  With
\[
  D_{\mathrm{ord}}(x)=x,\qquad
  L_{\mathrm{ord,full}}(x)=x,\qquad
  C=I_3,
\]
one has \(\XiC=0\).  If the ordinary-control coordinate
\(x_{\mathrm{ord}}\) is omitted, the trace residual is \(1\); the unit-root
control coordinate is therefore necessary for the full ordinary local
source.\footnote{Verified in Lean as
\texttt{SixBirdsBSD.Apparatus.PAdic.ordinarySchurCollapse}; see App~\ref{app:formalization}.}
\end{theorem}

\begin{proof}
With the full ordinary source admitted, both the readout and source maps are
the identity on \(\R^3\).  Hence
\[
  \KDD=I_3,\qquad
  \KLL=I_3,\qquad
  \Kdagger=I_3,\qquad
  \KDL=I_3,\qquad
  K_{LD}=I_3.
\]
The projection term is computed by the identity-matrix action.  For
\(v=(v_1,v_2,v_3)^{\mathsf T}\),
\[
  (I_3v)_i=\sum_{j=1}^3\delta_{ij}v_j=v_i,
\]
so \(I_3v=v\), and applying this column by column gives \(I_3I_3=I_3\).
Therefore
\[
  \KDL\Kdagger K_{LD}
  =
  I_3I_3I_3
  =
  I_3.
\]
Substituting into the Schur residual gives
\[
  \XiC
  =
  \KDD-\KDL\Kdagger K_{LD}
  =
  I_3-I_3
  =
  0.
\]

Now omit \(x_{\mathrm{ord}}\).  The admitted source is the coordinate
subspace spanned by \(e_f=(0,1,0)\) and \(e_{\exp}=(0,0,1)\).  The
orthogonal projector onto this subspace is
\[
  P_{\mathrm{omit}}
  =
  \operatorname{diag}(0,1,1).
\]
The residual projector is
\[
  I_3-P_{\mathrm{omit}}
  =
  \operatorname{diag}(1,0,0),
\]
whose trace is \(1\).  Thus omitting the ordinary-control coordinate leaves
exactly one unexplained local direction.
\end{proof}

\begin{theorem}[Signed local Schur collapse]\label{thm:apparatus:signed-schur-collapse}
For a good-supersingular row, let
\[
  y=(y_+,y_-,y_{\exp})\in\R^3
\]
be the signed \(+\), signed \(-\), and dual-exponential tangent coordinates.
With
\[
  D_{\pm}(y)=y,\qquad
  L_{\pm,\mathrm{full}}(y)=y,\qquad
  C=I_3,
\]
one has \(\XiC=0\).  The unsigned projection
\[
  L_{\mathrm{uns}}(y)=(y_++y_-,y_{\exp})
\]
has rank \(2\), leaves the difference direction \(y_+-y_-\) invisible, and
has trace residual \(1\).  Thus the signed decomposition is structural;
unsigned aggregation is a public shadow.\footnote{Verified in Lean as
\texttt{SixBirdsBSD.Apparatus.PAdic.signedSchurCollapse}; see App~\ref{app:formalization}.}
\end{theorem}

\begin{proof}
The full signed source is again the identity source on \(\R^3\).  Therefore
the same identity-matrix computation as in the ordinary case gives
\[
  \KDD=I_3,\qquad
  \KLL=I_3,\qquad
  \Kdagger=I_3,\qquad
  \KDL=I_3,\qquad
  K_{LD}=I_3,
\]
and hence
\[
  \XiC
  =
  I_3-I_3I_3I_3
  =
  0.
\]

It remains to compute the unsigned residual.  The unsigned source is the
two-dimensional subspace spanned by
\[
  u_1=(1,1,0),\qquad u_2=(0,0,1).
\]
These two vectors are orthogonal, with \(u_1\cdot u_1=2\),
\(u_2\cdot u_2=1\), and \(u_1\cdot u_2=0\).  Hence the orthogonal projector
onto the unsigned source is
\[
  P_{\mathrm{uns}}
  =
  \frac{u_1u_1^{\mathsf T}}{2}
  +
  u_2u_2^{\mathsf T}.
\]
Its trace is
\[
  \tr(P_{\mathrm{uns}})
  =
  \frac{u_1\cdot u_1}{2}
  +
  u_2\cdot u_2
  =
  \frac{2}{2}+1
  =
  2.
\]
Since \(\tr(I_3)=3\), the unsigned residual has trace
\[
  \tr(I_3-P_{\mathrm{uns}})=3-2=1.
\]

The invisible direction is explicit.  The vector
\[
  w=(1,-1,0)
\]
is orthogonal to \(u_1\) and \(u_2\), so it lies in the residual line.
Equivalently,
\[
  L_{\mathrm{uns}}(1,0,0)=(1,0)=L_{\mathrm{uns}}(0,1,0),
\]
although the two signed inputs are distinct.  Thus unsigned aggregation
forgets the \(+\) versus \(-\) difference, while the signed carrier retains
it.
\end{proof}

\begin{obligation}[Refined p-adic BK-local transport]\label{obl:apparatus:rho-p-transport}
The full Bloch--Kato local component requires a comparison transport
\[
  \beta_{\mathrm{BK},p}:
  \Delta_f(V,T)_{p,\R}
  \longrightarrow
  \Delta_f(M_E)_p.
\]
This transport fixes the integral Fontaine--Laffaille or Wach data, the
local Tate-pairing~\citep{Milne1986} and dual-exponential orientation, the Nekov\'a\v{r}
Selmer-complex local-global compatibility~\citep{Nekovar2006}, the signed normalization at
supersingular primes, and compatibility with the determinant-assembly
column.  In this paper it is a recognition-source carrier: the arithmetic
comparison is named external content and carried separately, not derived
from the finite typed local calculations above.
\end{obligation}

 \section{Determinant-assembly column and cascade completion}
\label{sec:det-assembly}

The determinant-assembly column records the orientation, sign, and unit
conventions needed to place the four arithmetic columns in one determinant
line.  The height/regulator, analytic/period, finite-source, and local
\(p\)-adic columns produce native determinant factors; the assembly column
is the final coordinate that makes their tensor a global determinant-line
object rather than a four-column product with an unspecified orientation.

\begin{definition}[Real determinant-assembly map]\label{def:apparatus:det-assembly-map}
Let
\[
  \lambda_{\mathrm{loc}}(E)
  =
  \rho_{\mathrm{ht/reg}}^{\R}
  \otimes
  \rho_{\mathrm{an/period}}^{\R}
  \otimes
  \rho_{\mathrm{finite}}^{\R}
  \otimes
  \bigotimes_{p\in S(E)}\rho_p^{\R}
\]
be the tensor of the four constructed arithmetic columns over the support
set \(S(E)\).  Let \(o_{\det}\) denote the determinant orientation and
trivialization coordinate.  The real determinant-assembly map is
\[
  \rho_{\det}^{\R}:
  \lambda_{\mathrm{loc}}(E)\otimes o_{\det}
  \longmapsto
  \lambda_{\mathrm{loc}}(E)\cdot o_{\det}
  \in\Delta_f(M_E)_{\R}.
\]
\end{definition}

\begin{theorem}[Cascade-completion Schur signature]\label{thm:apparatus:cascade-completion-signature}
Let the cascade tangent coordinate be
\[
  x=(x_{\mathrm{ht}},x_{\mathrm{an}},x_{\mathrm{fin}},x_p,x_{\det})\in\R^5.
\]
The four constructed arithmetic columns have source row
\[
  \ell=(1,1,1,1,0),
\]
while the global determinant target has row
\[
  d=(1,1,1,1,1).
\]
With \(C=I_5\), one has
\[
  \KDD=dd^{\mathsf T}=5,\qquad
  \KLL=\ell\ell^{\mathsf T}=4,\qquad
  \KDL=d\ell^{\mathsf T}=4,
\]
and therefore
\[
  \XiC=5-4\cdot4^{-1}\cdot4=1.
\]
The residual is exactly the determinant-orientation coordinate
\(x_{\det}\): the four constructed arithmetic components explain all
coordinates except the final determinant row.\footnote{Verified in Lean as
\texttt{SixBirdsBSD.Apparatus.DetAssembly.cascadeCompletionSignature}; see App~\ref{app:formalization}.}
\end{theorem}

\begin{proof}
The readout row is \(d=(1,1,1,1,1)\), so its squared norm is
\[
  \KDD
  =
  dd^{\mathsf T}
  =
  1+1+1+1+1
  =
  5.
\]
The source row supplied by the four constructed arithmetic columns is
\(\ell=(1,1,1,1,0)\).  Its squared norm is
\[
  \KLL
  =
  \ell\ell^{\mathsf T}
  =
  1+1+1+1+0
  =
  4.
\]
Since this is a one-dimensional source block, its Moore--Penrose inverse is
multiplication by \(4^{-1}\).  The cross term is
\[
  \KDL
  =
  d\ell^{\mathsf T}
  =
  1+1+1+1+0
  =
  4,
\]
and \(K_{LD}=4\) as well.  The Schur-complement residual is therefore
\[
  \XiC
  =
  \KDD-\KDL\Kdagger K_{LD}
  =
  5-4\cdot4^{-1}\cdot4
  =
  5-4
  =
  1.
\]

This scalar residual is precisely the determinant-orientation direction.
Indeed
\[
  d=(1,1,1,1,1)=\ell+e_{\det},
  \qquad
  e_{\det}=(0,0,0,0,1).
\]
The two summands are orthogonal for the identity calibration:
\(\ell\cdot e_{\det}=0\).  Projection onto the source row
\(\ell\) removes the \(\ell\)-component of \(d\) and leaves
\(e_{\det}\).  The squared norm of \(e_{\det}\) is \(1\), matching the
computed residual.  Thus the cascade is incomplete until the determinant
orientation coordinate is supplied.
\end{proof}

\begin{obligation}[Integral determinant assembly]\label{obl:apparatus:rho-det-assembly}
The full determinant comparison requires a transport
\[
  \beta_{\det\text{-}\mathrm{assembly}}:
  \Bigl(
  \Delta_{\mathrm{ht/reg}}
  \otimes
  \Delta_{\mathrm{an/period}}
  \otimes
  \Delta_{\mathrm{finite}}
  \otimes
  \bigotimes_p\Delta_p
  \Bigr)
  \longrightarrow
  \Delta_f(M_E).
\]
This transport fixes the Knudsen--Mumford determinant signs~\citep{KnudsenMumford1976}, the
Selmer-complex local-global compatibility, the zeta-element versus motivic
trivialization convention, and compatibility with the four refined
transports.  In this paper it is a recognition-source carrier: the
arithmetic determinant assembly is named external content and carried
separately, not derived from the finite cascade residual computation.
\end{obligation}

 \section[Comparison: literature gap and five-column/BK map]{Comparison: literature gap and $\Phi_{\mathrm{5col}\leftrightarrow\mathrm{BK}}$}
\label{sec:comparison}

The preceding sections built the five typed columns separately.  This
section contracts that typed atlas back to the usual Bloch--Kato scalar
package and asks whether the typed source is adequate for that contraction.
The answer is yes for the scalar readout: the typed source has zero Schur
residual for the comparison row, while the determinant-orientation
coordinate lies in the kernel of the scalar contraction.

\begin{remark}[Five-column literature gap (support only; audit observation)]\label{rmk:apparatus:five-column-literature-gap}
The literature audit used here reads Bloch--Kato~\citep[\S5]{BlochKato1990},
Fontaine--Perrin-Riou~\citep{FontainePerrinRiou1994}, and Nekov\'a\v{r}~\citep{Nekovar1994}
as presenting \(\Delta_f(M_E)\) as a single
global Tamagawa-number package, rather than as a typed decomposition into
height/regulator, analytic/period, finite-source, local \(p\)-adic, and
determinant-assembly slots.  Relative to that audited corpus, the
five-column decomposition is recorded here as a structural contribution of
this apparatus.  This is an audit observation about what that theorem-grade
literature supplies; it is not a non-existence claim about possible
equivalent decompositions elsewhere, and it is not a theorem of this paper.
\end{remark}

\begin{definition}[Five-column tensor]\label{def:apparatus:five-column-tensor}
The typed five-column tensor is
\[
  T_{\mathrm{5col}}(E)
  =
  \rho_{\mathrm{ht/reg}}^{\R}(E)
  \otimes
  \rho_{\mathrm{an/period}}^{\R}(E)
  \otimes
  \rho_{\mathrm{finite}}^{\R}(E)
  \otimes
  \bigotimes_{p\in S(E)}\rho_p^{\R}(E)
  \otimes
  \rho_{\det}^{\R}(E).
\]
Its source family \(L\) consists of the five separated log-tangent
coordinates
\[
  (x_{\mathrm{ht}},x_{\mathrm{an}},x_{\mathrm{fin}},x_p,x_{\det}),
\]
and its target readout \(D\) is the single Bloch--Kato
Tamagawa-package coordinate obtained after scalar contraction.
\end{definition}

\begin{theorem}[Five-column to BK scalar contraction]\label{thm:apparatus:five-col-bk-comparison}
At the log-tangent level, write
\[
  x=(x_{\mathrm{ht}},x_{\mathrm{an}},x_{\mathrm{fin}},x_p,x_{\det})\in\R^5.
\]
Define the five-column to Bloch--Kato scalar contraction by
\[
  \Phi_{\mathrm{5col}\leftrightarrow\mathrm{BK}}(x)=d\cdot x,
  \qquad
  d=(1,1,1,1,0).
\]
With the full typed five-column source \(L(x)=x\), \(C=I_5\),
\(\KLL=I_5\), \(\KDL=d\), and \(\KDD=\|d\|^2=4\), one has
\[
  \Xi_{\mathrm{5col}\leftrightarrow\mathrm{BK}}
  =
  4-dI_5d^{\mathsf T}
  =
  0.
\]
Thus the typed five-column source is adequate for the Bloch--Kato scalar
projection.  Moreover the determinant-orientation coordinate lies in the
kernel: \(x_{\det}\in\ker\Phi_{\mathrm{5col}\leftrightarrow\mathrm{BK}}\).
Consequently the scalar Bloch--Kato package is a strict shadow of the typed
atlas, forgetting at least the determinant-orientation coordinate.  The
integral arithmetic interpretation remains conditional on the five refined
transports.\footnote{Verified in Lean as
\texttt{SixBirdsBSD.Apparatus.Comparison.fiveColBkComparison}; see App~\ref{app:formalization}.}
\end{theorem}

\begin{proof}
The comparison row is \(d=(1,1,1,1,0)\).  Its squared norm is
\[
  \KDD
  =
  dd^{\mathsf T}
  =
  1+1+1+1+0
  =
  4.
\]
The full typed source is the identity source on \(\R^5\), so
\[
  \KLL=I_5,\qquad
  \Kdagger=I_5,\qquad
  \KDL=d,\qquad
  K_{LD}=d^{\mathsf T}.
\]
The projection term is again the identity-matrix action.  For any vector
\(v\in\R^5\), \(I_5v=v\), and hence
\[
  dI_5d^{\mathsf T}
  =
  dd^{\mathsf T}
  =
  4.
\]
Therefore
\[
  \Xi_{\mathrm{5col}\leftrightarrow\mathrm{BK}}
  =
  \KDD-\KDL\Kdagger K_{LD}
  =
  4-dI_5d^{\mathsf T}
  =
  4-4
  =
  0.
\]
This proves adequacy of the typed source for the scalar contraction.

It remains to identify the shadow.  Let
\[
  e_{\det}=(0,0,0,0,1).
\]
Then
\[
  d\cdot e_{\det}
  =
  1\cdot0+1\cdot0+1\cdot0+1\cdot0+0\cdot1
  =
  0.
\]
Thus \(e_{\det}\in\ker\Phi_{\mathrm{5col}\leftrightarrow\mathrm{BK}}\).
Changing only \(x_{\det}\) changes the typed five-column coordinate but
does not change the Bloch--Kato scalar readout.  The scalar package is
therefore a strict shadow of the typed atlas: it is adequately computed from
the typed source, but it does not remember all typed coordinates.
\end{proof}

 \section{Higher-rank no-go and the rank-2 Kummer source}
\label{sec:higher-rank-no-go}

The higher-rank no-go records a boundary of the height/regulator column.
Once a Mordell--Weil basis is known, the Neron--Tate Gram matrix computes
the regulator determinant.  It does not, by itself, construct independent
motivic-cycle sources in rank at least \(2\).  This is an audit result about
what the height matrix fails to determine; it is not a non-existence claim
about generalized Gross--Zagier~\citep{GrossZagier1986}, Beilinson--Kato~\citep{Beilinson1985,Kato2004}, \(K_2(E)\), or
higher-Chow constructions~\citep{Bloch1986} supplied by other means.

\begin{definition}[Rank-2 Mordell--Weil/Kummer source]\label{def:apparatus:rank-two-kummer-source}
Let
\[
  \Lambda_E=E(\Q)/E(\Q)_{\tors}.
\]
For a prime \(p\), the Kummer map
\[
  \kappa_p:\Lambda_E\otimes\Qp\longrightarrow \Hf(\Q,V_p(E))
\]
sends Mordell--Weil classes to finite Selmer classes.  For the curve
\(389a1\), fix a Mordell--Weil basis \((P_1,P_2)\).  The rank-\(2\)
Mordell--Weil/Kummer candidate is
\[
  c_2^{\mathrm{MW}}(P_i)=\kappa_p(P_i),\qquad i=1,2,
\]
viewed after tensoring with \(\Qp\).  This is a classical rank-\(2\)
Selmer source once the basis is known.  It is not, by this definition
alone, a generalized Heegner-cycle construction or a motivic-cycle source.
\end{definition}

\begin{theorem}[Rank-2 height Schur collapse]\label{thm:apparatus:rank-two-height-schur-collapse}
For the \(389a1\) rank-\(2\) Gram representative, write the symmetric
height variation in coordinates
\[
  (\delta h_{11},\delta h_{12},\delta h_{22})
\]
and write
\[
  d(\log\det)_H
  =
  \ell
  =
  (3.12679\ldots,-0.76771\ldots,2.14483\ldots).
\]
With the full symmetric height-matrix source, \(\KLL=I_3\) and
\[
  \KDD=\|\ell\|^2=14.96650\ldots.
\]
The Schur residual vanishes:
\[
  \Xi_{NT\to\Reg}
  =
  \|\ell\|^2-\ell I_3\ell^{\mathsf T}
  =
  0.
\]
Thus the full rank-\(2\) Neron--Tate Gram carrier is adequate for its
regulator determinant.  By contrast, the scalar-determinant-only source is
one-dimensional and leaves trace residual \(3-1=2\).\footnote{Verified in
Lean as
\texttt{SixBirdsBSD.Apparatus.HigherRankNoGo.rankTwoHeightSchurCollapse};
see App~\ref{app:formalization}.}
\end{theorem}

\begin{proof}
The tangent space of symmetric \(2\times2\) height matrices has the three
coordinates
\[
  (\delta h_{11},\delta h_{12},\delta h_{22}).
\]
The differential of the regulator determinant in logarithmic form is the
row vector \(\ell\).  Hence the readout energy before admitting a source is
\[
  \KDD
  =
  \ell\ell^{\mathsf T}
  =
  \|\ell\|^2
  =
  14.96650\ldots.
\]
When the full symmetric height-matrix carrier is admitted, the source block
is \(I_3\), its Moore--Penrose inverse is again \(I_3\), and the cross
blocks are \(\KDL=\ell\) and \(K_{LD}=\ell^{\mathsf T}\).  The projection
term is the identity-matrix action:
\[
  \ell I_3\ell^{\mathsf T}
  =
  \ell\ell^{\mathsf T}
  =
  \|\ell\|^2.
\]
Therefore
\[
  \Xi_{NT\to\Reg}
  =
  \KDD-\KDL\Kdagger K_{LD}
  =
  \|\ell\|^2-\|\ell\|^2
  =
  0.
\]

For the scalar-determinant-only comparison, the admitted source is the line
spanned by the nonzero vector \(\ell\) inside the three-dimensional
symmetric tangent space.  The orthogonal projector onto this line is
\[
  P_{\ell}=\frac{\ell^{\mathsf T}\ell}{\ell\ell^{\mathsf T}}.
\]
Its trace is
\[
  \tr(P_{\ell})
  =
  \frac{\ell\ell^{\mathsf T}}{\ell\ell^{\mathsf T}}
  =
  1.
\]
Since \(\tr(I_3)=3\), the scalar-determinant-only residual has trace
\[
  \tr(I_3-P_{\ell})=3-1=2.
\]
Thus the determinant scalar is adequate as a readout of the determinant
itself, but it does not retain the full three-coordinate height matrix.
\end{proof}

\begin{theorem}[Gram matrix is a noninjective shadow (NG-3)]\label{thm:apparatus:gram-matrix-shadow}
For \(r\geq2\), the Neron--Tate Gram matrix is a noninjective public shadow
of the rank-\(r\) motivic-cycle source.  The Gram construction records the
pairings
\[
  H_E=(\langle P_i,P_j\rangle_{\mathrm{NT}})_{1\leq i,j\leq r}
\]
for a known Mordell--Weil basis, and therefore supplies
\(\Reg(E/\Q)=\det H_E\).  It does not construct the independent rank-\(r\)
cycle source: distinct motivic-cycle data, including orientation, motivic
labels, and Abel--Jacobi provenance, can have the same Gram matrix.  Thus
the Mordell--Weil height matrix alone cannot be promoted to a rank-\(r\)
cycle-source carrier.\footnote{Verified in Lean as
\texttt{SixBirdsBSD.Apparatus.HigherRankNoGo.gramMatrixShadow}; see App~\ref{app:formalization}.}
\end{theorem}

\begin{proof}
Fix \(r\geq2\).  The Gram construction is a forgetful map.  Its input, if
one is seeking a motivic-cycle source, contains more than the numerical
height pairings: it must also specify which independent source classes are
being realized, their orientation, and their provenance through the relevant
cycle or Abel--Jacobi construction.  The Gram matrix keeps only the numbers
\(\langle P_i,P_j\rangle_{\mathrm{NT}}\).

This forgetfulness gives an explicit noninjectivity witness.  Take the same
Mordell--Weil basis \((P_1,\ldots,P_r)\) and attach two different motivic
source labels or two different Abel--Jacobi provenance records to the same
ordered basis.  These are distinct as motivic-cycle data, because the
source labels or provenance records differ.  Their Neron--Tate pairings are
identical, since the underlying Mordell--Weil points have not changed.
Therefore both data sets map to the same matrix
\[
  (\langle P_i,P_j\rangle_{\mathrm{NT}})_{i,j}.
\]
Equivalently, any height-orthogonal change of representatives that preserves
all pairings also leaves the Gram matrix unchanged while not supplying the
missing motivic labels or provenance.

Once a Mordell--Weil basis is known, the Gram matrix does compute the
regulator: \(\Reg(E/\Q)=\det H_E\).  That is exactly the content used by the
height/regulator column.  The additional claim needed for a rank-\(r\)
cycle source is different: it must construct or recognize independent
motivic classes whose regulator is being measured.  The Gram matrix contains
no such construction.  The no-go is therefore an audit result about the
naive route from height matrix to motivic source; it does not rule out BDP~\citep{BertoliniDarmonPrasanna2013},
Beilinson--Kato, \(K_2(E)\), or higher-Chow constructions imported under
their own hypotheses.
\end{proof}

The rank-\(\geq2\) generalized-Gross--Zagier gap is handled in the
conditional closure of \citet{TsiokosBSDClosure2026}, as a recognition
source rather than a closed theorem.

 \section{Global-audit no-go}
\label{sec:global-audit-no-go}

The global-audit no-go records the simplest obstruction to promoting a
finite table of elliptic curves to a global theorem.  The five-curve window
used in the audit is valuable evidence about those five rows, but a theorem
whose conclusion ranges over all curves over \(\Q\) cannot be determined
from those rows alone.  This is an audit result: it marks what the finite
window does not certify, and is not a claim that the missing global theorem
or a separate exhaustivity argument cannot exist.

\begin{theorem}[Finite-window promotion no-go (NG-4)]\label{thm:apparatus:finite-window-no-go}
Let
\[
  C_{\mathrm{BSD}}^{\mathrm{typed},+}
  =
  \operatorname*{colim}_{N_0\to\infty}
  \{E/\Q:\operatorname{cond}(E)\leq N_0\}
\]
be the conductor-colimit carrier of typed BSD rows, and let
\[
  W_{50}=\{11a1,37a1,121b1,960d1,571a1\}
\]
be the five-curve audit window.  No theorem-grade promotion from the
\(W_{50}\) rows to the global typed predicate
\[
  \forall E/\Q,\ \mathrm{BSD}(E)
\]
can depend only on \(W_{50}\).  In Schur form, at a conductor bound \(B\)
with \(M(B)=\# C_{\mathrm{BSD}}(B)\), the window source has
\(\KLL=I_5\), the full target has \(\KDD=I_{M(B)}\), and
\[
  \tr(\Xi_B)=M(B)-5\longrightarrow\infty
  \qquad (B\to\infty).
\]
Equivalently, two global completions agreeing on \(W_{50}\) but differing
at a curve \(E_*\notin W_{50}\) are indistinguishable to every
\(W_{50}\)-only operator.\footnote{Verified in Lean as
\texttt{SixBirdsBSD.Apparatus.GlobalAuditNoGo.finiteWindowNoGo}; see
App~\ref{app:formalization}.}
\end{theorem}

\begin{proof}
Fix a conductor bound \(B\) large enough that the bounded carrier
\[
  C_{\mathrm{BSD}}(B)=\{E/\Q:\operatorname{cond}(E)\leq B\}
\]
contains the five curves in \(W_{50}\).  Its typed audit space has one
coordinate for each row, so its target covariance block is the identity
matrix \(I_{M(B)}\), where \(M(B)=\#C_{\mathrm{BSD}}(B)\).  A source that
depends only on \(W_{50}\) sees exactly the five coordinate directions
corresponding to the five window rows.  Thus its source block is \(I_5\),
and the Schur projection is the coordinate projection from the
\(M(B)\)-dimensional row space onto those five window coordinates.

The trace of a coordinate projection is the number of retained coordinates.
Hence the window projection has trace \(5\), while the identity on the full
bounded carrier has trace \(M(B)\).  The unexplained residual is therefore
\[
  \tr(\Xi_B)
  =
  \tr(I_{M(B)})-\tr(P_{W_{50}})
  =
  M(B)-5.
\]
As the conductor bound increases, the carrier contains arbitrarily many
curves, so \(M(B)\) is unbounded and the residual trace tends to infinity.

The same obstruction can be seen without matrices.  Choose any
\(E_*\in C_{\mathrm{BSD}}^{\mathrm{typed},+}\setminus W_{50}\).  Consider two
formal global completions of the typed predicate.  They assign exactly the
same typed rows to all curves in \(W_{50}\), but assign different values at
\(E_*\); for instance, one completion marks the \(E_*\)-row as satisfying the
global predicate and the other marks it as not satisfying that predicate.
Any operator whose input is only the \(W_{50}\) table receives identical
data from these two completions, and therefore must return the same output
on both.  But the two completions differ on the global statement
\(\forall E/\Q,\mathrm{BSD}(E)\), since they differ at \(E_*\).  Therefore a
promotion depending only on the five window rows cannot determine the
global predicate.  A global closure requires some additional tail,
exhaustivity, or density-to-pointwise bridge; it is not supplied by the
finite window itself.
\end{proof}

 \section{Support-prime no-go and atlas}
\label{sec:support-prime-no-go}

The support-prime no-go is the local analogue of the finite-window
obstruction.  A finite list of primes can record useful local evidence, but
it cannot certify an all-prime support carrier by itself.  The result below
is therefore an audit result about non-exhaustive prime coverage, not a
claim that the missing local comparison rows cannot be supplied by other
arguments.

\begin{theorem}[Finite-prime cover no-go (NG-5)]\label{thm:apparatus:finite-prime-cover-no-go}
No fixed finite prime set certifies the all-prime support-prime carrier
\[
  \bigotimes_p \Delta_f(M_E)_p.
\]
More precisely, for every finite prime set \(P_0\), there is a prime
\(q\notin P_0\) and two support-prime completions which agree at every prime
in \(P_0\) but assign different coverage status to the \(q\)-local row.
Every \(P_0\)-only operator is therefore unable to distinguish the two
completions, and so cannot determine the all-prime target.\footnote{Verified
in Lean as
\texttt{SixBirdsBSD.Apparatus.SupportPrimeNoGo.finitePrimeCoverNoGo}; see
App~\ref{app:formalization}.}
\end{theorem}

\begin{proof}
Let \(P_0\) be any finite set of primes.  Since the set of primes is
infinite, choose a prime \(q\notin P_0\).  We now compare two formal
support-prime completions for the same global row.  They agree on all local
data indexed by primes in \(P_0\).  At the new prime \(q\), the first
completion marks the local row as covered, while the second marks it as
uncovered.

By definition, a \(P_0\)-only operator factors through the restriction of
the all-prime support carrier to the primes in \(P_0\).  The two completions
have the same restriction to \(P_0\), so such an operator receives identical
input from them and must return the same value on both.  But the all-prime
target distinguishes them, because the \(q\)-coordinate has different
coverage status.  Hence the \(P_0\)-restricted data do not determine the
all-prime carrier.

Since \(P_0\) was arbitrary, no fixed finite list of primes certifies
\(\bigotimes_p\Delta_f(M_E)_p\).  Finite-prime evidence can support a
bounded audit, but all-prime coverage requires an additional exhaustivity or
local-completion certificate for the residual support-prime rows.
\end{proof}

\begin{definition}[Support-prime truncation residual (finite diagnostic)]\label{def:apparatus:support-prime-truncation-residual}
For a finite prime truncation
\[
  S_B=\{p:p\leq B\},
\]
classify each local row into the covered class R4a or the residual class
R4b.  Here R4a denotes a row whose ordinary or signed local hypotheses
contribute no residual in the audit, while R4b denotes a row with one unit
of unexplained local data.  The finite-truncation residual is
\[
  \Xi_B(E)=u_B(E),
\]
where \(u_B(E)\) is the number of R4b rows in \(S_B\).  This is a finite
diagnostic only, distinct from
\Cref{thm:apparatus:finite-prime-cover-no-go}; it does not by itself prove
all-prime non-exhaustivity.  For the b50 rows below \(50\),
\[
  \Xi(11a1)=\Xi(37a1)=\Xi(121b1)=1,\qquad
  \Xi(960d1)=3,\qquad
  \Xi(571a1)=0.
\]
The last value becomes \(1\) once the bad prime \(571\) is included.
\end{definition}

\begin{remark}[Support-prime atlas (support only)]\label{rmk:apparatus:support-prime-atlas}
The support-prime atlas is a \(76\)-row audit table of pairs \((E,p)\) for
the b50 curves.  For each row it records the reduction type, the value
\(a_p\) obtained by point counting, ordinarity, applicability of the local
Iwasawa main-conjecture input~\citep{SkinnerUrban2014}, and the R4a/R4b assignment.  The R4b rows
below \(50\) are
\[
  11a1:\{11\},\qquad
  37a1:\{37\},\qquad
  121b1:\{11\},\qquad
  960d1:\{2,3,5\},\qquad
  571a1:\varnothing.
\]
The prime \(571\) for \(571a1\) lies outside the \(<50\) sample, and the
b50 bad-prime R4b count is \(7\).  This atlas is a support table for the
audit, not a theorem.
\end{remark}

The full apparatus no-go suite is collected in \Cref{tab:no-go-suite}.  The
shared reading is that each row marks a loss of source information in the
available public or finite record, not a theorem excluding future richer
source data.
\begin{table}[tbp]
\centering
\footnotesize
\caption{The apparatus no-go suite.  Each item is an audit result about what
the displayed public row does not determine; none is a non-existence theorem
about richer external source data.}
\label{tab:no-go-suite}
\begin{tabularx}{\textwidth}{@{}>{\raggedright\arraybackslash}p{0.15\textwidth}
>{\raggedright\arraybackslash}p{0.27\textwidth}
>{\raggedright\arraybackslash}p{0.36\textwidth}
>{\raggedright\arraybackslash}X@{}}
\toprule
No-go & Form & What the literature or public row lacks & Status \\
\midrule
NG-1 & Finite scalar shadow & The scalar finite factor does not identify
which finite-source block carries the valuation. & mechanized finite
noninjectivity witness. \\
\addlinespace
NG-2 & \(\dim \Sha[p]\) shadow & Dimension-only data forgets the
Cassels--Tate determinant and pairing structure. & mechanized projector
trace residual. \\
\addlinespace
NG-3 & Gram-matrix/regulator shadow & A regulator determinant does not
construct the rank-\(2\) source whose Gram matrix it would measure. &
mechanized typed shadow statement. \\
\addlinespace
NG-4 & Finite-window promotion & A bounded audit window cannot certify an
unbounded operator family. & mechanized finite-window obstruction. \\
\addlinespace
NG-5 & Finite-prime coverage & A fixed finite prime set cannot determine an
all-prime support carrier. & mechanized finite-prime cover obstruction. \\
\addlinespace
Higher-rank GZ & Rank-\(\geq2\) source gap & The theorem-grade
generalized Gross--Zagier identity is not supplied by the audited inputs. &
recorded as an audit-result gap. \\
\addlinespace
Global audit & Operator promotion gap & Local or finite scalar evidence does
not by itself promote to the global determinant-line carrier. & recorded as
scope boundary. \\
\addlinespace
Support-prime atlas & Local support gap & The \(b50\) support table is finite
evidence, not all-prime coverage. & support-only audit table. \\
\bottomrule
\end{tabularx}
\end{table}

 \section{Vsrc: Non-Descending Translation calibration}
\label{sec:vsrc}

The \(\vsrc\) calibration isolates the determinant-shaped source content
that is visible before any descent to actual motivic cycles is supplied.
In rank \(2\) it is the exterior algebra on two formal source symbols:
the product \(\Jone\Jtwo\) records independence, while the relations
\(\Jone^2=\Jtwo^2=0\) and \(\Jtwo\Jone=-\Jone\Jtwo\) make the construction
alternating rather than Clifford-like.

\begin{definition}[Virtual source algebra]\label{def:apparatus:vsrc-algebra}
For an elliptic curve \(E/\Q\) whose source problem is considered in rank
\(2\), let
\[
  V_2=\Q\Jone\oplus\Q\Jtwo.
\]
The virtual source algebra is
\[
  \vsrc_2(E)=\Lambda_{\Q}(V_2),
\]
with unit \(1\) and ordered basis
\[
  1,\qquad \Jone,\qquad \Jtwo,\qquad \Jtwelve:=\Jone\Jtwo.
\]
Its multiplication is the exterior product, so
\[
  \Jone^2=0,\qquad \Jtwo^2=0,\qquad
  \Jone\Jtwo=\Jtwelve,\qquad
  \Jtwo\Jone=-\Jtwelve.
\]
The symbols \(\Jone\) and \(\Jtwo\) are non-descending virtual independent
source classes; the element \(\Jtwelve\) is their determinant-shaped
independence witness.\footnote{Realized in Lean as
\texttt{SixBirdsBSD.Apparatus.Vsrc.vsrcAlgebra} over a narrowed
representation: the rank-2 exterior-algebra relations are encoded
structurally on the \(2^r\) coordinate space (dimension/calibration content
faithful; full axiomatization not claimed). See App~\ref{app:formalization}.}
\end{definition}

\Cref{fig:vsrc-wedge} displays the rank-\(2\) exterior product used in the
definition.  The diagram is only the finite structural calibration; it is not a
descent to motivic cycles.
\begin{figure}[tbp]
\centering
\begin{tikzpicture}[
  font=\small,
  node/.style={draw, rounded corners=2pt, align=center, minimum width=1.6cm,
    minimum height=0.8cm, fill=black!3},
  arrow/.style={-{Latex[length=2mm]}, thick}
]
\node[node] (one) at (0,2.2) {\(1\)};
\node[node] (j1) at (-2.2,0.7) {\(\Jone\)};
\node[node] (j2) at (2.2,0.7) {\(\Jtwo\)};
\node[node] (j12) at (0,-1.1) {\(\Jtwelve\)};

\draw[arrow] (j1) -- node[above, sloped] {\(\wedge\,\Jtwo\)} (j12);
\draw[arrow] (j2) -- node[above, sloped] {\(\wedge\,\Jone=-\)} (j12);
\draw[arrow] (one) -- node[left] {\(\wedge\,\Jone\)} (j1);
\draw[arrow] (one) -- node[right] {\(\wedge\,\Jtwo\)} (j2);
\draw[loop left, looseness=6, -{Latex[length=2mm]}] (j1) to node[left]
  {\(\Jone^2=0\)} (j1);
\draw[loop right, looseness=6, -{Latex[length=2mm]}] (j2) to node[right]
  {\(\Jtwo^2=0\)} (j2);
\node[align=center, font=\footnotesize] at (0,-2.25)
  {\(\dim_{\Q}\vsrc_2=4\), with \(\Jone\Jtwo=\Jtwelve\) the determinant-shaped witness};
\end{tikzpicture}
\caption{The rank-\(2\) \(\vsrc\) exterior-algebra calibration.}
\label{fig:vsrc-wedge}
\end{figure}

\begin{theorem}[Stage I consistency]\label{thm:apparatus:vsrc-stage-i-consistency}
The rank-\(2\) virtual source algebra is a consistent finite-dimensional
associative unital graded-commutative \(\Q\)-algebra.  Every element has a
unique ordered normal form
\[
  a_{\varnothing}\,1+a_1\Jone+a_2\Jtwo+a_{12}\Jtwelve,
  \qquad a_{\varnothing},a_1,a_2,a_{12}\in\Q,
\]
and therefore
\[
  \dim_{\Q}\vsrc_2(E)=4=2^2.
\]
The low-rank cases are preserved: \(\vsrc_0(E)=\Q\), and
\(\vsrc_1(E)=\Q\oplus \Q J_1\) with \(J_1^2=0\).  Under the usual
Gross--Zagier hypotheses, the rank-\(1\) symbol is the one that later
descends to the Heegner point \(P_K\).\footnote{Verified in Lean as
\texttt{SixBirdsBSD.Apparatus.Vsrc.vsrcStageIConsistency} over the narrowed
\(2^r\)-coordinate representation (structural encoding; full axiomatization
not claimed). See App~\ref{app:formalization}.}
\end{theorem}

\begin{proof}
Construct the algebra on the \(\Q\)-vector space with basis indexed by the
subsets of \(\{1,2\}\).  The empty subset gives the basis vector \(1\), the
singletons give \(\Jone\) and \(\Jtwo\), and the full subset gives the
ordered product \(\Jtwelve=\Jone\Jtwo\).  On basis elements define the
product as follows.  If two subsets meet, their product is \(0\).  If they
are disjoint, their product is the basis element of the union, multiplied
by the sign needed to move the concatenated list into increasing order.
Extend this rule \(\Q\)-bilinearly.

This rule immediately gives the displayed rank-\(2\) multiplication table:
\[
  \Jone^2=0,\qquad \Jtwo^2=0,
\]
because each singleton meets itself, and
\[
  \Jone\Jtwo=\Jtwelve,\qquad
  \Jtwo\Jone=-\Jtwelve,
\]
because reversing the two singleton factors introduces one transposition.
The empty subset is a two-sided unit, since adjoining the empty list changes
neither the union nor the sign.

Associativity may be checked on basis vectors and then extended by
bilinearity.  If any index is repeated among three basis factors, both
parenthesizations are zero: at the step where the repeated index is
multiplied into a factor already containing it, the intersection is
nonempty.  If all indices appearing in the three factors are distinct, then
both parenthesizations produce the basis vector of the same union.  The
sign in each case is the parity of the same permutation that orders the
three concatenated lists; splitting the multiplication into two stages only
splits the inversion count into two summands.  Thus the signs agree, and
the product is associative.

The graded-commutativity is the same sign rule: moving a homogeneous basis
monomial of degree \(a\) past one of degree \(b\) contributes
\((-1)^{ab}\), and an overlap gives zero on both sides.  Hence the algebra
is associative, unital, and graded-commutative.

Finally, the four subset-indexed basis vectors are linearly independent by
construction, so every element has the unique normal form
\[
  a_{\varnothing}\,1+a_1\Jone+a_2\Jtwo+a_{12}\Jtwelve.
\]
There are \(2^2=4\) subsets of \(\{1,2\}\), giving dimension \(4\).  The
same subset construction in rank \(0\) has only the empty subset and gives
\(\Q\); in rank \(1\) it has the two subsets \(\varnothing\) and
\(\{1\}\), giving \(\Q\oplus\Q J_1\) with \(J_1^2=0\).  This proves the
Stage-I consistency assertions.
\end{proof}

\begin{remark}[No licensed rank-$r$ descent (support only; ledger predicate)]\label{rmk:apparatus:vsrc-no-licensed-descent}
For \(r\geq2\), the inherited audit records do not license a descent
\[
  \vsrc_r(E)\longrightarrow \mathrm{CH}^*(\mathcal X)_{\Q}
\]
or a descent to a Selmer or Mordell--Weil carrier that realizes all symbols
\(J_i\) as independent motivic classes.  This is a statement about the
recorded inputs: the Gram-matrix no-go leaves motivic labels and provenance
unconstructed; the Kato Euler-system~\citep{Kato2004} audit does not supply \(r\) independent
source classes; and the BDP~\citep{BertoliniDarmonPrasanna2013} generalized-Heegner audit records motivic
provenance without a two-class closure for \(389a1\).  The remark is a
support-only ledger observation, not a mathematical non-existence claim
about possible future or externally supplied descents.
\end{remark}

\begin{theorem}[Stage II Gross--Zagier partial descent]\label{thm:apparatus:vsrc-stage-ii-gz-partial}
In rank \(1\), the \(\vsrc_1\) calculus recovers the Gross--Zagier~\citep{GrossZagier1986,Kolyvagin1990}
source/regulator side under the usual Heegner hypotheses.  If
\[
  r_{\mathrm{GZ}}(J_1)=P_K
\]
and the virtual height form is sent to the Neron--Tate height pairing, then
\[
  h_{\vsrc}(J_1,J_1)
  \longmapsto
  \langle P_K,P_K\rangle_{\mathrm{NT}},
  \qquad
  \Reg=\langle P_K,P_K\rangle_{\mathrm{NT}}.
\]
The same symbolic descent does not recover the full Gross--Zagier identity
\[
  \frac{L'(E/K,1)}{\Omega_{E,K}}
  =
  c(E,K)\,\langle P_K,P_K\rangle_{\mathrm{NT}}
\]
unless the analytic-height theorem is supplied separately.  Thus the strict
descent residual is the named analytic-height import
\[
  \Xi_{\mathrm{GZ}}
  =
  E_{\mathrm{analytic\text{-}height\text{-}import}}
  \neq 0
\]
relative to the source/regulator calculus alone.\footnote{Verified in Lean
as \texttt{SixBirdsBSD.Apparatus.Vsrc.vsrcStageIIGzPartial}; see App~\ref{app:formalization}.}
\end{theorem}

\begin{proof}
The rank-\(1\) algebra is
\[
  \vsrc_1(E)=\Q\oplus\Q J_1,\qquad J_1^2=0.
\]
Under the Gross--Zagier hypotheses one has a Heegner point \(P_K\).  Define
the rank-\(1\) descent on the source generator by
\[
  r_{\mathrm{GZ}}(J_1)=P_K
\]
and send the virtual height form to the Neron--Tate height form.  Then the
unique nontrivial regulator entry in rank \(1\) is carried to
\[
  h_{\vsrc}(J_1,J_1)
  \mapsto
  \langle P_K,P_K\rangle_{\mathrm{NT}}.
\]
For a rank-\(1\) Mordell--Weil lattice, the regulator is exactly the
\(1\times1\) determinant of this height matrix, hence
\[
  \Reg=\langle P_K,P_K\rangle_{\mathrm{NT}}.
\]
This proves the source/regulator part.

The full Gross--Zagier formula contains additional analytic information:
the leading derivative \(L'(E/K,1)\), the period \(\Omega_{E,K}\), and the
normalizing factor \(c(E,K)\).  These quantities do not appear in the
rank-\(1\) exterior algebra or in the descent of \(J_1\) to \(P_K\).  Indeed,
two formal data sets can have the same \(P_K\) and the same
\(\langle P_K,P_K\rangle_{\mathrm{NT}}\) while assigning different analytic
leading quotients \(L'(E/K,1)/\Omega_{E,K}\).  The symbolic descent therefore
cannot determine the displayed analytic-height identity.  The missing
assertion is precisely the external analytic-height theorem, recorded here
as the residual
\[
  E_{\mathrm{analytic\text{-}height\text{-}import}}.
\]
Thus Stage II recovers the Gross--Zagier source/regulator side, but strict
descent to the full identity requires that separate theorem.
\end{proof}

\begin{theorem}[Stage III translation (389a1)]\label{thm:apparatus:vsrc-stage-iii-translation}
For \(389a1\), take the rank-\(2\) algebra
\[
  \vsrc_2(E)=\Q\oplus\Q\Jone\oplus\Q\Jtwo\oplus\Q\Jtwelve,
  \qquad
  \Jone\Jtwo=\Jtwelve.
\]
Since the finite factors for this calibration satisfy
\[
  \prod c_v=\#\Sha=|E_{\tors}|=1,
\]
the scalar residual is
\[
  \Xi_{\mathrm{Sc}}=0
  \quad\Longleftrightarrow\quad
  A_{\mathrm{lead}}=\Omega_E^+\,\Reg(E/\Q).
\]
The motivic residual \(\Xi_{\mathrm{Mot}}\) vanishes exactly when there is a
rank-\(2\) realization
\[
  (r_2,z_{12},\mathcal B_{\mathrm{an\text{-}ht}}^{(2)})
\]
with
\[
  z_{12}=r_2(\Jtwelve)\neq0
\]
and with \(\mathcal B_{\mathrm{an\text{-}ht}}^{(2)}\) identifying the
rank-\(2\) analytic leading term with the height/regulator determinant of
that wedge source.  Under this bridge,
\[
  \Xi_{\mathrm{Mot}}=0\quad\Longrightarrow\quad \Xi_{\mathrm{Sc}}=0.
\]
The reverse implication is not established by the inherited records: the
scalar equality does not construct \(z_{12}\), does not certify its
independence, and therefore requires the named external content
\(\mathcal R_{\mathrm{BSD}\to\mathrm{mot}}^{(2)}\).\footnote{Verified in
Lean as \texttt{SixBirdsBSD.Apparatus.Vsrc.vsrcStageIIITranslation}; see
App~\ref{app:formalization}.}
\end{theorem}

\begin{proof}
Assume first that \(\Xi_{\mathrm{Mot}}=0\).  By the stated meaning of the
motivic residual, this gives a realization \(r_2\), a nonzero wedge source
\[
  z_{12}=r_2(\Jtwelve),
\]
and a rank-\(2\) analytic-height bridge
\(\mathcal B_{\mathrm{an\text{-}ht}}^{(2)}\).  The bridge identifies the
analytic leading coefficient attached to the rank-\(2\) source with the
height determinant of that source.  For the fixed \(389a1\) calibration, the
finite correction factors are all \(1\):
\[
  \prod c_v=\#\Sha=|E_{\tors}|=1.
\]
Thus the bridge gives exactly
\[
  A_{\mathrm{lead}}=\Omega_E^+\,\Reg(E/\Q),
\]
which is the condition \(\Xi_{\mathrm{Sc}}=0\).  This proves the forward
translation.

Now suppose only that \(\Xi_{\mathrm{Sc}}=0\).  This is the scalar equality
\[
  A_{\mathrm{lead}}=\Omega_E^+\,\Reg(E/\Q).
\]
It is an equality of numerical determinant readouts.  It contains no choice
of realization map \(r_2\), no element \(z_{12}\), and no certificate that a
nonzero image of \(\Jtwelve\) is an independent motivic wedge source.  The
Gram-matrix shadow obstruction from the higher-rank no-go is exactly this
loss of source data: the regulator can be known without constructing the
rank-\(2\) motivic source whose regulator it would measure.  Therefore the
reverse implication is not supplied by the scalar equality.  It would
require the separate recognition content
\[
  \mathcal R_{\mathrm{BSD}\to\mathrm{mot}}^{(2)}
\]
together with the corresponding bridge and independence certificate.  The
statement is consequently a one-way translation, not a converse failure.
\end{proof}

\begin{remark}[Rank-$r$ stronger-than-scalar reading (support only)]\label{rmk:apparatus:vsrc-rank-r-stronger}
The one-way Stage-III translation exposes a structural asymmetry: the scalar
leading-coefficient formula records a determinant value, while the rank-\(r\)
source statement also asks for the existence and independence of the wedge
class
\[
  J_1\wedge\cdots\wedge J_r.
\]
In that sense, the rank-\(r\) source reading carries more information than
the scalar equality.  This is a support-only interpretation of the
Stage-III asymmetry, not an independently proven meta-theorem that scalar
BSD is strictly weaker in every external theory.
\end{remark}

 \section[Canonical kappa-r normalization]{$\kapr$ normalization}
\label{sec:kappa-normalization}

The normalization constant \(\kapr\) is the scalar that makes the cascade
leading-term form use the same arithmetic factor as the standard Strong BSD
leading-coefficient formula.  The point of this section is only to identify
that scalar after the cascade convention \(c_E=\Tam(E)\) is fixed.  It is an
equivalence of two ways of writing the same conjectural identity, not a
derivation of BSD.

\begin{theorem}[Canonical $\kapr$ normalization equivalence]\label{thm:apparatus:kappa-r-normalization-equivalence}
Fix the cascade convention
\[
  c_E:=\Tam(E)=\prod_{\ell} c_{\ell}(E).
\]
Then the cascade leading-term form
\[
  \frac{L^{(r)}(E,1)}{r!}
  =
  \kapr(E)\,\Omega_E\,\Reg_r(E)\,c_E
\]
agrees with the standard Strong BSD leading-coefficient form
\[
  \frac{L^{(r)}(E,1)}{r!}
  =
  \Omega_E\,\Reg_r(E)\,
  \frac{\#\Sha(E/\Q)\,\Tam(E)}{\#E(\Q)_{\tors}^{\,2}}
\]
if and only if
\[
  \kapr(E)=
  \frac{\#\Sha(E/\Q)}{\#E(\Q)_{\tors}^{\,2}}.
\]
Under this canonical normalization, the cascade form is equivalent to the
standard Strong BSD identity as a restatement in cascade target coordinates.
It does not prove the identity.\footnote{\raggedright Verified in Lean as
\texttt{SixBirdsBSD.\allowbreak Apparatus.\allowbreak KappaNormalization.\allowbreak kappaRNormalizationEquivalence};
see App~\ref{app:formalization}.}
\end{theorem}

\begin{proof}
After the convention \(c_E=\Tam(E)\), the right-hand side of the cascade
form is
\[
  \kapr(E)\,\Omega_E\,\Reg_r(E)\,\Tam(E).
\]
The right-hand side of the standard Strong BSD leading-coefficient formula
is
\[
  \Omega_E\,\Reg_r(E)\,
  \frac{\#\Sha(E/\Q)\,\Tam(E)}{\#E(\Q)_{\tors}^{\,2}}
  =
  \frac{\#\Sha(E/\Q)}{\#E(\Q)_{\tors}^{\,2}}\,
  \Omega_E\,\Reg_r(E)\,\Tam(E).
\]
Thus the two formulas have the same left-hand side and differ only in the
coefficient multiplying the common factor
\[
  \Omega_E\,\Reg_r(E)\,\Tam(E).
\]
In the BSD normalization this factor is nonzero: the period is positive,
the regulator determinant is positive in positive rank and equal to \(1\)
in rank \(0\), and the Tamagawa product is a positive integer.  Therefore
the two right-hand sides are equal if and only if their scalar coefficients
are equal, namely if and only if
\[
  \kapr(E)
  =
  \frac{\#\Sha(E/\Q)}{\#E(\Q)_{\tors}^{\,2}}.
\]

If this equality of coefficients is imposed, substituting it into the
cascade form gives exactly the standard Strong BSD leading-coefficient
formula.  Conversely, if the cascade form and the standard form agree under
the fixed convention \(c_E=\Tam(E)\), cancellation of the same nonzero common
factor forces the displayed value of \(\kapr(E)\).  The proof is only this
algebraic comparison of normalizations; it supplies no proof that either
leading-coefficient identity is true for \(E\).
\end{proof}

The conditional Strong-BSD closure of \citet{TsiokosBSDClosure2026}
consumes this normalization.  The role of the present section is only to
supply the normalization equivalence used there.
 \section{Scope and non-claims}\label{sec:scope}


The preceding sections build a typed apparatus for comparing the usual
BSD/Bloch--Kato scalar packages with finer source carriers.  The apparatus
is deliberately limited.  It separates what is proved at the finite typed
level from what is supplied by external arithmetic input, and it records
several audit no-gos without turning them into mathematical non-existence
claims.

The four non-claims governing this paper are the following.

\begin{itemize}
  \item \textbf{A-NC-1 --- No-gos are audit results, not non-existence claims.}
  The no-go suite (NG-1--NG-5, higher-rank, global, support-prime)
  records where theorem-grade Mode-A literature does not reach; it is not
  a claim that the missing identities cannot exist.  Here ``Mode-A''
  denotes the theorem-grade arithmetic literature audited as external
  provenance, as opposed to a finite diagnostic, numerical example, or
  support table.

  \item \textbf{A-NC-2 --- Finite typed level.}
  Adequacy collapses are proved at the finite, mathlib-free, typed level
  (identity-matrix action; \(\Q\) rank-1 projector traces);
  operator/motivic full generality is not claimed.  Thus the Schur
  residual calculations in this paper are finite vector and matrix
  computations over explicit typed carriers.  They are not asserted as
  analytic operator theorems or as constructions of motivic cycles.

  \item \textbf{A-NC-3 --- \(\vsrc\) structural narrowing.}
  The \(\vsrc\) exterior-product relations are represented structurally on
  the \(2^r\) coordinate space; the dimension/calibration content is
  faithful, full axiomatization is not claimed.  The \(\vsrc\) section
  therefore calibrates the determinant
  shape of a rank-\(r\) source and the low-rank relations used in the
  paper; it does not present a complete general exterior-algebra
  formalization or a descent to actual motivic cycles.

  \item \textbf{A-NC-4 --- Recognition-source transports are carriers.}
  The five per-column transports are recognition sources (carriers), not
  derived theorems.  They are supplied or imported as named arithmetic
  comparison content and carried as hypotheses where needed; they are not
  established from the finite typed apparatus itself.
\end{itemize}

Consequently, this paper should not be read as proving BSD, an Iwasawa main
conjecture, or a Gross--Zagier theorem.  The column computations prove
adequacy or residual statements for the typed carriers introduced here.
The recognition sources and imports are supplied, imported, or carried as
hypotheses; they are never claimed to be established by the apparatus.
 \section{Discussion}\label{sec:discussion}


The apparatus developed here is meant to be reusable.  Its basic question is
not whether a given BSD or Bloch--Kato identity is true, but which typed
source data are adequate for a specified readout.  The Schur residual
\(\XiC\) makes that question finite and local to a chosen column: if the
native source coordinates act as the identity on the tangent readout, the
residual vanishes; if a public scalar keeps only a projection, the residual
measures the lost directions.  This turns a familiar informal observation
about scalar invariants into a computable diagnostic.  One asks, column by
column, whether the available datum is the source itself or only its public
shadow.

This diagnostic is useful because the standard BSD formula is naturally
read as a product of scalar factors.  The leading analytic coefficient is
divided by a period; the height pairing contributes its determinant; finite
arithmetic enters through Tamagawa numbers, torsion, and \(\Sha\); local
\(p\)-adic data are typically compressed through ordinary, signed, or
control-theoretic normalizations; and determinant signs and units are
hidden in the final assembly.  The scalar formula is therefore a powerful
summary, but it is not a complete description of the source data behind
each factor.  The five-column decomposition refines this summary into
height/regulator, analytic/period, finite-source, \(p\)-adic, and
determinant-assembly columns.  It keeps the sources separated until the
comparison map deliberately contracts them back to the usual
Bloch--Kato/Tamagawa-number package.

Seen from the classical Bloch--Kato perspective, the point is not to replace
the fundamental line by a new object.  The point is to expose the internal
typing that is suppressed when the global Tamagawa-number package is treated
as a single scalar row.  The comparison theorem shows that the typed
five-column source is adequate for the Bloch--Kato scalar contraction: the
Schur residual for that contraction is zero.  At the same time, the kernel
contains the determinant-orientation coordinate.  Thus the usual scalar
package is recovered, but it is recovered as a projection of a richer typed
atlas.  This is exactly what one expects from a determinant-line
formalism: the global scalar is the visible readout, while the determinant
trivializations, local orientations, and source-specific normalizations
explain how the readout is assembled.

The per-column collapses have a similarly modest interpretation.  The
height/regulator column collapses because the full Neron--Tate pairing is
admitted as source data; the analytic/period column collapses because both
\(\log A_E\) and \(\log\Omega_E^+\) are present; the finite-source column
collapses only when Tamagawa, torsion, \(\Sha[p^\infty]\), and
Cassels--Tate pairing data are all retained; the \(p\)-adic columns collapse
only after including the ordinary-control coordinate or separating the
signed supersingular directions.  These are finite identity-matrix
computations.  Their value is that they make explicit which coordinates
must be present before a row deserves to be treated as a source rather than
as a quotient or shadow.

The no-go suite records the complementary boundary.  The finite scalar
factor forgets where a valuation lies: a contribution in the \(\Sha\)-block
and a contribution in the Tamagawa block can have the same scalar value.
The dimension-only \(\Sha[p]\) readout forgets the Cassels--Tate
determinant data.  The Neron--Tate Gram matrix gives the regulator once a
Mordell--Weil basis is known, but it does not construct independent
motivic-cycle sources.  A finite curve window cannot certify a universal
statement depending on all curves, and a finite prime list cannot certify
all local support-prime rows.  None of these no-gos says that the missing
arithmetic identities or constructions do not exist.  They say that the
particular public row under inspection is not enough to determine the
typed target without additional source data or an external recognition
input.

Recording that boundary precisely matters.  In classical prose, it is easy
to slide from ``this scalar agrees numerically'' to ``this source has been
constructed,'' or from ``this finite sample supports the pattern'' to
``this global statement is certified.''  The Schur residual prevents that
slide by asking what the operator can see.  If two typed completions agree
on the public row but differ in the target coordinate, then the public row
is not an adequate source for that target.  Conversely, when the full typed
source is admitted and the residual vanishes by an identity action, the
adequacy statement is exactly as strong as the declared source.  This
separation is the central conceptual point of the apparatus.

The \(\vsrc\) calibration guards the source end of the story.  In rank
\(2\), the exterior product \(\Jone\Jtwo=\Jtwelve\) is a compact way to
record determinant-shaped independence before any descent to motivic
cycles has been supplied.  The Stage I algebra is consistent, and the
Stage II rank-\(1\) descent recovers the Gross--Zagier source/regulator
side once a Heegner point is available.  Stage III then explains the
asymmetry in rank \(2\): a motivic wedge source together with the analytic
height bridge implies the scalar leading-coefficient equality, while the
scalar equality alone does not construct the wedge source.  This is the
same lesson as the Gram-matrix shadow no-go, now expressed at the source
end of the apparatus.

The \(\kapr\) normalization guards the target end.  Once
\[
  c_E=\Tam(E)
\]
is fixed, the cascade leading-term form agrees with the standard Strong
BSD leading-coefficient form exactly when
\[
  \kapr(E)=\frac{\#\Sha(E/\Q)}{\#E(\Q)_{\tors}^{\,2}}.
\]
This is only an equality of normalizations, but it is an important one:
it prevents the cascade coordinates from silently changing the arithmetic
constant.  Together, the \(\vsrc\) calibration and the \(\kapr\)
normalization keep the apparatus honest at both ends.  The source side
does not pretend that scalar regulator data construct motivic classes, and
the target side does not pretend that a differently normalized cascade
formula is the standard Strong BSD formula.

The framework lineage supplies the language of typed carriers, visibility,
and residual accounting \citep{TsiokosNeedles2026,TsiokosFoundationsII2026,
TsiokosFoundationsIII2026}.  The mathematical content of this paper is the
specialization of that language to the BSD/Bloch--Kato setting: a
five-column typed atlas, finite Schur adequacy calculations, audit no-gos
for public shadows, and normalization checks at the source and target
interfaces.  The resulting layer can be consumed by later conditional
closures or by independent audits, but it is already useful on its own: it
turns the question ``what does this BSD scalar know?'' into a precise
finite calculation about what the chosen source row retains and what it
forgets.
 \section{Conclusion}\label{sec:conclusion}


This paper has isolated a typed five-column decomposition of the
BSD/Bloch--Kato fundamental-line data and used finite Schur-complement
calculations to test adequacy column by column.  The height/regulator,
analytic/period, finite-source, \(p\)-adic, and determinant-assembly
columns each come with an explicit residual calculation; the comparison
map contracts the resulting typed atlas back to the standard scalar
Bloch--Kato/Tamagawa package.  The accompanying no-go suite (NG-1--NG-5,
higher-rank, global, and support-prime) marks where a public scalar row
underdetermines its typed target: scalar finite factors forget source
blocks, dimension-only \(\Sha[p]\) data forget Cassels--Tate structure,
Gram matrices do not construct motivic sources, finite curve windows do
not certify global statements, and finite prime lists do not certify
all-prime support.

The boundary is part of the result.  The adequacy collapses proved here
are finite, typed, identity-action statements; they do not assert full
operator or motivic generality.  The recognition-source transports and
external arithmetic identities are carried as hypotheses or imported
comparison content, not proved by the apparatus.  The open gaps left
explicit include an all-prime exhaustivity certificate, a
density-to-pointwise or global-promotion bridge, and a genuinely general
operator/motivic form of the column collapses.

The reusable value is the diagnostic itself: the apparatus turns the
question ``what does this scalar know?'' into a finite, per-column,
checkable question about which source directions are retained and which
are lost.
 \appendix
\section{Lean formalization}\label{app:formalization}
\subsection*{Trust base and methodology}

The Lean development for this paper is mathlib-free.  It imports no
external mathematics library, so no mathlib axioms enter the dependency
closure.  Every declaration registered as a theorem was audited with
\texttt{\#print axioms}.  The only axioms that occur are Lean's standard
foundational ones: \texttt{propext}, \texttt{Quot.sound}, and
\texttt{Classical.choice}, with \texttt{funext} and the \texttt{choice}
recursor appearing as their standard consequences in the audited closures.
There are no project-local axioms behind any theorem-marked result.  Thus
the trust base for the formalized theorem claims is precisely Lean's kernel
plus classical logic, and nothing project-specific.

This limited trust base is possible because the formalized content is
finite and typed.  The adequacy collapses and no-go residuals are not
mechanized as analytic estimates, operator-theoretic assertions, or
motivic constructions.  They are finite vector and matrix computations:
the source and target covariance blocks are explicit finite matrices, the
full-source collapses reduce to identity-matrix action, and the scalar
shadow residuals are computed from rank-\(1\) projectors over \(\Q\).
Each Schur-complement residual is therefore a decidable calculation in a
small finite algebraic model.  This is the reason a mathlib-free kernel can
check the main adequacy and no-go statements.

The surrounding arithmetic geometry is treated separately from those finite
calculations.  The Lean development checks the typed identities appearing
in the apparatus: for example, that an identity source kills the Schur
residual, that the relevant projector traces have the stated values, and
that the finite noninjectivity witnesses have the claimed public shadows.
It does not certify Gross--Zagier, Iwasawa main conjecture inputs,
Cassels--Tate comparison theory, Bloch--Kato local comparison maps, or the
motivic descents named in the body.  Such arithmetic content enters the
paper either as cited external input, as a recognition-source carrier, or
as an explicitly stated hypothesis.

The statements-of-record inventory for this paper contains \(41\) items.
Of these, \(19\) are theorem-grade Lean results with substantive finite
proofs, \(11\) are typed definition mirrors, and \(11\) are not mechanized:
the five recognition-source transport obligations, the four support-only
remarks, the warning, and the scope nonclaim.  Those non-mechanized items
are paper-level statements carried as motivation, scope, or hypotheses,
not Lean declarations.  Two mechanized items, the \(\vsrc\) algebra
definition and the Stage-I \(\vsrc\) consistency theorem, are checked over
a narrowed structural representation; the following subsections record the
per-label coverage, declaration map, and standing narrowings in detail.
\subsection*{Wording-discipline table}

Each theorem in the body carries exactly one mechanization footnote, and
the wording of that footnote is canonical rather than editorial.  The
wording records the fidelity class of the formal counterpart: a faithful
Lean theorem, a Lean theorem over a narrowed representation, a typed
definition mirror, or an unmechanized paper-level statement.  Narrowing and
non-mechanization are therefore visible in the disclosure convention rather
than hidden in prose.

\begin{center}
\begin{tabularx}{\textwidth}{@{}>{\raggedright\arraybackslash}p{0.31\textwidth}X@{}}
\toprule
Statement class & Footnote convention \\
\midrule
Theorem, faithful &
Verified in Lean as \texttt{[decl]}; see App~\ref{app:formalization}. \\
\addlinespace
Theorem, narrowed representation &
Verified in Lean as \texttt{[decl]} over a narrowed representation
(e.g.\ the \(2^r\)-coordinate space); the narrowing is recorded below.  The
narrowing qualifier is mandatory; never bare ``verified in Lean''. \\
\addlinespace
Definition, narrowed representation &
Realized in Lean as \texttt{[decl]} over a narrowed representation; the
narrowing is recorded below. \\
\addlinespace
Definition, faithful (typed mirror) &
No footnote; the declaration appears only in the per-label coverage table
below. \\
\addlinespace
Obligation / remark / warning / non-claim (not mechanized) &
No footnote and no Lean identifier in the body; recorded as a paper-level
statement only. \\
\bottomrule
\end{tabularx}
\end{center}
\subsection*{Standing narrowings}

There are two standing qualifications behind the coverage claims in this
appendix.  Both are part of the formalization contract: they delimit what
the finite Lean development checks, and they keep the surrounding
arithmetic content visible rather than implicit.

First, two \(\vsrc\) items are checked over a narrowed structural
representation: the \(\vsrc\) algebra definition and the Stage-I
\(\vsrc\) consistency theorem.  On the \(2^r\)-coordinate model, the
determinant-shaped content of a rank-\(r\) source is represented, and the
low-rank exterior relations actually used in the paper are checked.  In
rank \(2\), this includes the wedge relation
\[
  \Jone\Jtwo=\Jtwelve
\]
together with the accompanying nilpotence and anticommutation relations.
The dimension count and the calibration of the determinant shape are
faithful to the intended exterior-algebra source.

What is not claimed is equally important.  The formalization is not a
complete general exterior-algebra development: it does not provide general
\(\bigwedge^k\) functoriality, nor a full axiomatization of exterior
algebras in arbitrary rank.  It is also not a descent from \(\vsrc\) to
actual motivic cycles.  The Lean check certifies the finite structural
relations on the coordinate representation, not a general motivic
construction.  This is why those two body items carry the ``narrowed
representation'' footnote convention rather than a bare ``verified in
Lean'' disclosure, and why the \(\vsrc\) statements in
\Cref{sec:vsrc} are calibration and consistency results rather than
constructions of motivic sources.

Second, the five per-column transports are recognition sources, not
derived theorems and not Lean declarations.  These are the
height/regulator, analytic/period, finite-source, \(p\)-adic, and
determinant-assembly comparison transports appearing in
\Cref{sec:height-regulator,sec:analytic-period,sec:finite-source,sec:p-adic,sec:det-assembly}.
Each column's Schur-adequacy collapse is proved at the finite typed level
once the column's native source is admitted.  The corresponding transport
is precisely the named arithmetic comparison content that admits that
source into the Bloch--Kato determinant-line target.  It is supplied or
imported externally and then carried as a hypothesis.

Thus the apparatus does not prove these five transports.  It records them
as obligations and proves the per-column collapses conditional on the
native source data that those transports are meant to recognize.  In the
per-label coverage table below they appear as not-mechanized obligations.
Neither standing narrowing enlarges the Lean trust base: no new axioms are
introduced; the qualifications bound the finite mechanization claims, with
the remaining arithmetic content carried explicitly.
\begin{table}[tbp]
\centering
\small
\caption{Standing representation notes for the apparatus formalization.}
\label{tab:representation-notes}
\begin{tabularx}{\textwidth}{@{}>{\raggedright\arraybackslash}p{0.24\textwidth}
>{\raggedright\arraybackslash}p{0.28\textwidth}
>{\raggedright\arraybackslash}X@{}}
\toprule
Representation issue & Where it appears & Disclosure effect \\
\midrule
\(\vsrc\) structural encoding &
\(\Cref{sec:vsrc}\), App~\ref{app:formalization} &
The rank-\(2\) exterior relations are checked on the \(2^r\)-coordinate
model.  The paper claims calibration and consistency, not a full exterior
algebra or motivic descent formalization. \\
\addlinespace
\(\Q\) rank-\(1\) projector traces &
Finite-source and no-go statements &
Scalar-shadow residuals are finite rational projector computations; the
Lean result is the finite trace calculation, not an arithmetic
classification theorem. \\
\addlinespace
Identity-matrix action &
Schur-collapse statements &
Full-source adequacy reduces to \(I\)-action on an explicit finite carrier.
The external transport that admits the native source remains a carried
recognition obligation. \\
\bottomrule
\end{tabularx}
\end{table}
 \subsection*{Per-label coverage}

The following table lists every statement of record in document order,
with its environment kind, its mechanization fidelity, and its body
location.  The paper labels in the first column elide the
\texttt{apparatus:} prefix.  Faithful and narrowed rows have Lean
declarations, mapped in the next subsection; not-mechanized rows do not.

\begin{small}
\begin{longtable}{@{}>{\ttfamily}p{0.45\textwidth} l l l@{}}
\caption{Per-label coverage for the apparatus statements of record (41 rows, document order). The first column is the paper label (the \texttt{apparatus:} prefix elided); the last column points to the body location.}\label{tab:formalization-coverage}\\
\toprule
\normalfont Paper label & Kind & Fidelity & Body \\
\midrule
\endfirsthead
\multicolumn{4}{@{}l}{\footnotesize\itshape continued from previous page}\\
\toprule \normalfont Paper label & Kind & Fidelity & Body \\ \midrule
\endhead
\midrule \multicolumn{4}{r@{}}{\footnotesize\itshape continued on next page}\\
\endfoot
\bottomrule
\endlastfoot
def:bk-component-map & \normalfont definition & \normalfont faithful & \normalfont \Cref{sec:decomposition} \\
def:bridge-defect-equation & \normalfont definition & \normalfont faithful & \normalfont \Cref{sec:decomposition} \\
thm:component-killing & \normalfont theorem & \normalfont faithful & \normalfont \Cref{sec:decomposition} \\
thm:finite-scalar-public-shadow & \normalfont theorem & \normalfont faithful & \normalfont \Cref{sec:decomposition} \\
nonclaim:decomposition-scope & \normalfont nonclaim & \normalfont not mechanized & \normalfont \Cref{sec:decomposition} \\
def:height-regulator-map & \normalfont definition & \normalfont faithful & \normalfont \Cref{sec:height-regulator} \\
thm:height-schur-collapse & \normalfont theorem & \normalfont faithful & \normalfont \Cref{sec:height-regulator} \\
obl:rho-ht-reg-transport & \normalfont obligation & \normalfont not mechanized & \normalfont \Cref{sec:height-regulator} \\
warn:571a1-carrier & \normalfont warning & \normalfont not mechanized & \normalfont \Cref{sec:height-regulator} \\
def:analytic-period-map & \normalfont definition & \normalfont faithful & \normalfont \Cref{sec:analytic-period} \\
thm:analytic-period-schur-collapse & \normalfont theorem & \normalfont faithful & \normalfont \Cref{sec:analytic-period} \\
obl:rho-an-period-transport & \normalfont obligation & \normalfont not mechanized & \normalfont \Cref{sec:analytic-period} \\
def:finite-source-map & \normalfont definition & \normalfont faithful & \normalfont \Cref{sec:finite-source} \\
thm:finite-source-schur-collapse & \normalfont theorem & \normalfont faithful & \normalfont \Cref{sec:finite-source} \\
thm:scalar-finite-shadow & \normalfont theorem & \normalfont faithful & \normalfont \Cref{sec:finite-source} \\
thm:dim-sha-shadow & \normalfont theorem & \normalfont faithful & \normalfont \Cref{sec:finite-source} \\
obl:rho-finite-transport & \normalfont obligation & \normalfont not mechanized & \normalfont \Cref{sec:finite-source} \\
def:p-adic-map & \normalfont definition & \normalfont faithful & \normalfont \Cref{sec:p-adic} \\
thm:ordinary-schur-collapse & \normalfont theorem & \normalfont faithful & \normalfont \Cref{sec:p-adic} \\
thm:signed-schur-collapse & \normalfont theorem & \normalfont faithful & \normalfont \Cref{sec:p-adic} \\
obl:rho-p-transport & \normalfont obligation & \normalfont not mechanized & \normalfont \Cref{sec:p-adic} \\
def:det-assembly-map & \normalfont definition & \normalfont faithful & \normalfont \Cref{sec:det-assembly} \\
thm:cascade-completion-signature & \normalfont theorem & \normalfont faithful & \normalfont \Cref{sec:det-assembly} \\
obl:rho-det-assembly & \normalfont obligation & \normalfont not mechanized & \normalfont \Cref{sec:det-assembly} \\
rmk:five-column-literature-gap & \normalfont remark & \normalfont not mechanized & \normalfont \Cref{sec:comparison} \\
def:five-column-tensor & \normalfont definition & \normalfont faithful & \normalfont \Cref{sec:comparison} \\
thm:five-col-bk-comparison & \normalfont theorem & \normalfont faithful & \normalfont \Cref{sec:comparison} \\
def:rank-two-kummer-source & \normalfont definition & \normalfont faithful & \normalfont \Cref{sec:higher-rank-no-go} \\
thm:rank-two-height-schur-collapse & \normalfont theorem & \normalfont faithful & \normalfont \Cref{sec:higher-rank-no-go} \\
thm:gram-matrix-shadow & \normalfont theorem & \normalfont faithful & \normalfont \Cref{sec:higher-rank-no-go} \\
thm:finite-window-no-go & \normalfont theorem & \normalfont faithful & \normalfont \Cref{sec:global-audit-no-go} \\
thm:finite-prime-cover-no-go & \normalfont theorem & \normalfont faithful & \normalfont \Cref{sec:support-prime-no-go} \\
def:support-prime-truncation-residual & \normalfont definition & \normalfont faithful & \normalfont \Cref{sec:support-prime-no-go} \\
rmk:support-prime-atlas & \normalfont remark & \normalfont not mechanized & \normalfont \Cref{sec:support-prime-no-go} \\
def:vsrc-algebra & \normalfont definition & \normalfont narrowed & \normalfont \Cref{sec:vsrc} \\
thm:vsrc-stage-i-consistency & \normalfont theorem & \normalfont narrowed & \normalfont \Cref{sec:vsrc} \\
rmk:vsrc-no-licensed-descent & \normalfont remark & \normalfont not mechanized & \normalfont \Cref{sec:vsrc} \\
thm:vsrc-stage-ii-gz-partial & \normalfont theorem & \normalfont faithful & \normalfont \Cref{sec:vsrc} \\
thm:vsrc-stage-iii-translation & \normalfont theorem & \normalfont faithful & \normalfont \Cref{sec:vsrc} \\
rmk:vsrc-rank-r-stronger & \normalfont remark & \normalfont not mechanized & \normalfont \Cref{sec:vsrc} \\
thm:kappa-r-normalization-equivalence & \normalfont theorem & \normalfont faithful & \normalfont \Cref{sec:kappa-normalization} \\
\end{longtable}
\end{small}
 \subsection*{Lean declaration map}

The \(30\) mechanized statements, consisting of the \(11\) typed
definitions and \(19\) theorems, each correspond to a Lean declaration
listed here for citation and auditing.  Paper labels elide the
\texttt{apparatus:} prefix, and declarations elide the
\texttt{SixBirdsBSD.Apparatus.} namespace.  The not-mechanized
obligations, remarks, warning, and nonclaim have no declaration and are
therefore omitted; the full coverage table, including those rows, is
\Cref{tab:formalization-coverage}.

\begin{small}
\begin{longtable}{@{}>{\ttfamily\raggedright\arraybackslash}p{0.44\textwidth} >{\ttfamily\raggedright\arraybackslash}p{0.50\textwidth}@{}}
\caption{Lean declaration map for the 30 mechanized apparatus statements (definitions and theorems). Paper labels elide the \texttt{apparatus:} prefix; declarations elide the \texttt{SixBirdsBSD.Apparatus.} namespace.}\label{tab:appd-declmap}\\
\toprule \normalfont Paper label & \normalfont Lean declaration \\ \midrule
\endfirsthead
\multicolumn{2}{@{}l}{\footnotesize\itshape continued from previous page}\\
\toprule \normalfont Paper label & \normalfont Lean declaration \\ \midrule
\endhead
\midrule \multicolumn{2}{r@{}}{\footnotesize\itshape continued on next page}\\
\endfoot
\bottomrule
\endlastfoot
def:\allowbreak bk-\allowbreak component-\allowbreak map & Decomposition.\allowbreak bk\allowbreak Component\allowbreak Map \\
def:\allowbreak bridge-\allowbreak defect-\allowbreak equation & Decomposition.\allowbreak bridge\allowbreak Defect\allowbreak Equation \\
thm:\allowbreak component-\allowbreak killing & Decomposition.\allowbreak component\allowbreak Killing \\
thm:\allowbreak finite-\allowbreak scalar-\allowbreak public-\allowbreak shadow & Decomposition.\allowbreak finite\allowbreak Scalar\allowbreak Public\allowbreak Shadow \\
def:\allowbreak height-\allowbreak regulator-\allowbreak map & Height\allowbreak Regulator.\allowbreak height\allowbreak Regulator\allowbreak Map \\
thm:\allowbreak height-\allowbreak schur-\allowbreak collapse & Height\allowbreak Regulator.\allowbreak height\allowbreak Schur\allowbreak Collapse \\
def:\allowbreak analytic-\allowbreak period-\allowbreak map & Analytic\allowbreak Period.\allowbreak analytic\allowbreak Period\allowbreak Map \\
thm:\allowbreak analytic-\allowbreak period-\allowbreak schur-\allowbreak collapse & Analytic\allowbreak Period.\allowbreak analytic\allowbreak Period\allowbreak Schur\allowbreak Collapse \\
def:\allowbreak finite-\allowbreak source-\allowbreak map & Finite\allowbreak Source.\allowbreak finite\allowbreak Source\allowbreak Map \\
thm:\allowbreak finite-\allowbreak source-\allowbreak schur-\allowbreak collapse & Finite\allowbreak Source.\allowbreak finite\allowbreak Source\allowbreak Schur\allowbreak Collapse \\
thm:\allowbreak scalar-\allowbreak finite-\allowbreak shadow & Finite\allowbreak Source.\allowbreak scalar\allowbreak Finite\allowbreak Shadow \\
thm:\allowbreak dim-\allowbreak sha-\allowbreak shadow & Finite\allowbreak Source.\allowbreak dim\allowbreak Sha\allowbreak Shadow \\
def:\allowbreak p-\allowbreak adic-\allowbreak map & PAdic.\allowbreak p\allowbreak Adic\allowbreak Map \\
thm:\allowbreak ordinary-\allowbreak schur-\allowbreak collapse & PAdic.\allowbreak ordinary\allowbreak Schur\allowbreak Collapse \\
thm:\allowbreak signed-\allowbreak schur-\allowbreak collapse & PAdic.\allowbreak signed\allowbreak Schur\allowbreak Collapse \\
def:\allowbreak det-\allowbreak assembly-\allowbreak map & Det\allowbreak Assembly.\allowbreak det\allowbreak Assembly\allowbreak Map \\
thm:\allowbreak cascade-\allowbreak completion-\allowbreak signature & Det\allowbreak Assembly.\allowbreak cascade\allowbreak Completion\allowbreak Signature \\
def:\allowbreak five-\allowbreak column-\allowbreak tensor & Comparison.\allowbreak five\allowbreak Column\allowbreak Tensor \\
thm:\allowbreak five-\allowbreak col-\allowbreak bk-\allowbreak comparison & Comparison.\allowbreak five\allowbreak Col\allowbreak Bk\allowbreak Comparison \\
def:\allowbreak rank-\allowbreak two-\allowbreak kummer-\allowbreak source & Higher\allowbreak Rank\allowbreak No\allowbreak Go.\allowbreak rank\allowbreak Two\allowbreak Kummer\allowbreak Source \\
thm:\allowbreak rank-\allowbreak two-\allowbreak height-\allowbreak schur-\allowbreak collapse & Higher\allowbreak Rank\allowbreak No\allowbreak Go.\allowbreak rank\allowbreak Two\allowbreak Height\allowbreak Schur\allowbreak Collapse \\
thm:\allowbreak gram-\allowbreak matrix-\allowbreak shadow & Higher\allowbreak Rank\allowbreak No\allowbreak Go.\allowbreak gram\allowbreak Matrix\allowbreak Shadow \\
thm:\allowbreak finite-\allowbreak window-\allowbreak no-\allowbreak go & Global\allowbreak Audit\allowbreak No\allowbreak Go.\allowbreak finite\allowbreak Window\allowbreak No\allowbreak Go \\
thm:\allowbreak finite-\allowbreak prime-\allowbreak cover-\allowbreak no-\allowbreak go & Support\allowbreak Prime\allowbreak No\allowbreak Go.\allowbreak finite\allowbreak Prime\allowbreak Cover\allowbreak No\allowbreak Go \\
def:\allowbreak support-\allowbreak prime-\allowbreak truncation-\allowbreak residual & Support\allowbreak Prime\allowbreak No\allowbreak Go.\allowbreak support\allowbreak Prime\allowbreak Truncation\allowbreak Residual \\
def:\allowbreak vsrc-\allowbreak algebra & Vsrc.\allowbreak vsrc\allowbreak Algebra \\
thm:\allowbreak vsrc-\allowbreak stage-\allowbreak i-\allowbreak consistency & Vsrc.\allowbreak vsrc\allowbreak Stage\allowbreak IConsistency \\
thm:\allowbreak vsrc-\allowbreak stage-\allowbreak ii-\allowbreak gz-\allowbreak partial & Vsrc.\allowbreak vsrc\allowbreak Stage\allowbreak IIGz\allowbreak Partial \\
thm:\allowbreak vsrc-\allowbreak stage-\allowbreak iii-\allowbreak translation & Vsrc.\allowbreak vsrc\allowbreak Stage\allowbreak IIITranslation \\
thm:\allowbreak kappa-\allowbreak r-\allowbreak normalization-\allowbreak equivalence & Kappa\allowbreak Normalization.\allowbreak kappa\allowbreak RNormalization\allowbreak Equivalence \\
\end{longtable}
\end{small}

\clearpage
\begingroup
\raggedright
\small
\bibliographystyle{plainnat}
\begin{thebibliography}{37}
\providecommand{\natexlab}[1]{#1}
\providecommand{\url}[1]{\texttt{#1}}
\expandafter\ifx\csname urlstyle\endcsname\relax
  \providecommand{\doi}[1]{doi: #1}\else
  \providecommand{\doi}{doi: \begingroup \urlstyle{rm}\Url}\fi

\bibitem[Beilinson(1985)]{Beilinson1985}
A.~A. Beilinson.
\newblock Higher regulators and values of {L}-functions.
\newblock \emph{J. Soviet Math.}, 30:\penalty0 2036--2070, 1985.
\newblock \doi{10.1007/BF02105861}.

\bibitem[Bertolini et~al.(2013)Bertolini, Darmon, and Prasanna]{BertoliniDarmonPrasanna2013}
Massimo Bertolini, Henri Darmon, and Kartik Prasanna.
\newblock Generalized {Heegner} cycles and p-adic {Rankin} {L}-series.
\newblock \emph{Duke Math. J.}, 162\penalty0 (6):\penalty0 1033--1148, 2013.
\newblock \doi{10.1215/00127094-2142056}.

\bibitem[Birch and Swinnerton-Dyer(1965)]{BirchSwinnertonDyer1965}
B.~J. Birch and H.~P.~F. Swinnerton-Dyer.
\newblock Notes on elliptic curves. {II}.
\newblock \emph{J. Reine Angew. Math.}, 218:\penalty0 79--108, 1965.
\newblock \doi{10.1515/crll.1965.218.79}.

\bibitem[Bloch(1986)]{Bloch1986}
Spencer Bloch.
\newblock Algebraic cycles and higher {K}-theory.
\newblock \emph{Adv. Math.}, 61\penalty0 (3):\penalty0 267--304, 1986.
\newblock \doi{10.1016/0001-8708(86)90081-2}.

\bibitem[Bloch and Kato(1990)]{BlochKato1990}
Spencer Bloch and Kazuya Kato.
\newblock {L}-functions and {T}amagawa numbers of motives.
\newblock In \emph{The {Grothendieck} {Festschrift}, Vol.~I}, volume~86 of \emph{Progr. Math.}, pages 333--400. Birkh\"auser, 1990.
\newblock \doi{10.1007/978-0-8176-4574-8_9}.

\bibitem[Breuil et~al.(2001)Breuil, Conrad, Diamond, and Taylor]{BreuilConradDiamondTaylor2001}
Christophe Breuil, Brian Conrad, Fred Diamond, and Richard Taylor.
\newblock On the modularity of elliptic curves over {$\mathbf{Q}$}: wild 3-adic exercises.
\newblock \emph{J. Amer. Math. Soc.}, 14\penalty0 (4):\penalty0 843--939, 2001.
\newblock \doi{10.1090/S0894-0347-01-00370-8}.

\bibitem[Burns and Flach(2001)]{BurnsFlach2001}
David Burns and Matthias Flach.
\newblock Tamagawa numbers for motives with (non-commutative) coefficients.
\newblock \emph{Doc. Math.}, 6:\penalty0 501--570, 2001.
\newblock \doi{10.4171/DM/113}.

\bibitem[Cassels(1962)]{Cassels1962}
J.~W.~S. Cassels.
\newblock Arithmetic on curves of genus 1. {IV}. {Proof} of the {Hauptvermutung}.
\newblock \emph{J. Reine Angew. Math.}, 211:\penalty0 95--112, 1962.

\bibitem[Cremona(1997)]{Cremona1997}
J.~E. Cremona.
\newblock \emph{Algorithms for Modular Elliptic Curves}.
\newblock Cambridge Univ. Press, 2 edition, 1997.

\bibitem[Deligne(1987)]{Deligne1987}
Pierre Deligne.
\newblock Le d\'eterminant de la cohomologie.
\newblock In \emph{Current Trends in Arithmetical Algebraic Geometry}, volume~67 of \emph{Contemp. Math.}, pages 93--177. Amer. Math. Soc., 1987.

\bibitem[Fontaine(1994)]{Fontaine1994}
Jean-Marc Fontaine.
\newblock Expos\'e {II}: Le corps des p\'eriodes {$p$}-adiques.
\newblock In Jean-Marc Fontaine, editor, \emph{P\'eriodes {$p$}-adiques -- S\'eminaire de {Bures}, 1988}, number 223 in Ast\'erisque, pages 59--101. Soc. Math. France, 1994.
\newblock URL \url{https://www.numdam.org/item/AST_1994__223__59_0/}.

\bibitem[Fontaine and Perrin-Riou(1994)]{FontainePerrinRiou1994}
Jean-Marc Fontaine and Bernadette Perrin-Riou.
\newblock Autour des conjectures de {Bloch} et {Kato}: cohomologie galoisienne et valeurs de fonctions {L}.
\newblock In \emph{Motives ({Seattle}, 1991)}, volume 55.1 of \emph{Proc. Sympos. Pure Math.}, pages 599--706. Amer. Math. Soc., 1994.

\bibitem[Gross and Zagier(1986)]{GrossZagier1986}
Benedict~H. Gross and Don~B. Zagier.
\newblock Heegner points and derivatives of {L}-series.
\newblock \emph{Invent. Math.}, 84\penalty0 (2):\penalty0 225--320, 1986.
\newblock \doi{10.1007/BF01388809}.

\bibitem[Horn and Johnson(2013)]{HornJohnson2013}
Roger~A. Horn and Charles~R. Johnson.
\newblock \emph{Matrix Analysis}.
\newblock Cambridge Univ. Press, 2 edition, 2013.

\bibitem[Kato(2004)]{Kato2004}
Kazuya Kato.
\newblock p-adic {Hodge} theory and values of zeta functions of modular forms.
\newblock \emph{Ast\'erisque}, 295:\penalty0 117--290, 2004.

\bibitem[Knudsen and Mumford(1976)]{KnudsenMumford1976}
Finn~Faye Knudsen and David Mumford.
\newblock The projectivity of the moduli space of stable curves. {I}. {Preliminaries} on ``det'' and ``div''.
\newblock \emph{Math. Scand.}, 39\penalty0 (1):\penalty0 19--55, 1976.
\newblock \doi{10.7146/math.scand.a-11642}.

\bibitem[Kobayashi(2003)]{Kobayashi2003}
Shin-ichi Kobayashi.
\newblock Iwasawa theory for elliptic curves at supersingular primes.
\newblock \emph{Invent. Math.}, 152\penalty0 (1):\penalty0 1--36, 2003.
\newblock \doi{10.1007/s00222-002-0265-4}.

\bibitem[Kolyvagin(1990)]{Kolyvagin1990}
V.~A. Kolyvagin.
\newblock Euler systems.
\newblock In \emph{The {Grothendieck} {Festschrift}, Vol.~II}, volume~87 of \emph{Progr. Math.}, pages 435--483. Birkh\"auser, 1990.

\bibitem[Mazur(1972)]{Mazur1972}
Barry Mazur.
\newblock Rational points of abelian varieties with values in towers of number fields.
\newblock \emph{Invent. Math.}, 18:\penalty0 183--266, 1972.
\newblock \doi{10.1007/BF01389815}.

\bibitem[Milne(1986)]{Milne1986}
J.~S. Milne.
\newblock \emph{Arithmetic Duality Theorems}, volume~1 of \emph{Perspectives in Mathematics}.
\newblock Academic Press, 1986.

\bibitem[Nekov\'a\v{r}(1994)]{Nekovar1994}
Jan Nekov\'a\v{r}.
\newblock {Beilinson's} conjectures.
\newblock In \emph{Motives ({Seattle}, 1991)}, volume 55.1 of \emph{Proc. Sympos. Pure Math.}, pages 537--570. Amer. Math. Soc., 1994.

\bibitem[Nekov\'a\v{r}(2006)]{Nekovar2006}
Jan Nekov\'a\v{r}.
\newblock \emph{Selmer Complexes}.
\newblock Number 310 in Ast\'erisque. Soc. Math. France, 2006.

\bibitem[N\'eron(1965)]{Neron1965}
Andr\'e N\'eron.
\newblock Quasi-fonctions et hauteurs sur les vari\'et\'es ab\'eliennes.
\newblock \emph{Ann. of Math. (2)}, 82:\penalty0 249--331, 1965.
\newblock \doi{10.2307/1970644}.

\bibitem[Pollack(2003)]{Pollack2003}
Robert Pollack.
\newblock On the p-adic {L}-function of a modular form at a supersingular prime.
\newblock \emph{Duke Math. J.}, 118\penalty0 (3):\penalty0 523--558, 2003.
\newblock \doi{10.1215/S0012-7094-03-11835-9}.

\bibitem[Poonen and Stoll(1999)]{PoonenStoll1999}
Bjorn Poonen and Michael Stoll.
\newblock The {Cassels--Tate} pairing on polarized abelian varieties.
\newblock \emph{Ann. of Math. (2)}, 150\penalty0 (3):\penalty0 1109--1149, 1999.
\newblock \doi{10.2307/121064}.

\bibitem[Schneider(1982)]{Schneider1982}
Peter Schneider.
\newblock p-adic height pairings. {I}.
\newblock \emph{Invent. Math.}, 69\penalty0 (3):\penalty0 401--409, 1982.
\newblock \doi{10.1007/BF01389362}.

\bibitem[Silverman(1994)]{Silverman1994}
Joseph~H. Silverman.
\newblock \emph{Advanced Topics in the Arithmetic of Elliptic Curves}, volume 151 of \emph{Graduate Texts in Mathematics}.
\newblock Springer, 1994.
\newblock \doi{10.1007/978-1-4612-0851-8}.

\bibitem[Silverman(2009)]{Silverman2009}
Joseph~H. Silverman.
\newblock \emph{The Arithmetic of Elliptic Curves}, volume 106 of \emph{Graduate Texts in Mathematics}.
\newblock Springer, 2 edition, 2009.
\newblock \doi{10.1007/978-0-387-09494-6}.

\bibitem[Skinner and Urban(2014)]{SkinnerUrban2014}
Christopher Skinner and Eric Urban.
\newblock The {Iwasawa} main conjectures for {$\mathrm{GL}_2$}.
\newblock \emph{Invent. Math.}, 195\penalty0 (1):\penalty0 1--277, 2014.
\newblock \doi{10.1007/s00222-013-0448-1}.

\bibitem[Tate(1966)]{Tate1966bsd}
John~T. Tate.
\newblock On the conjectures of {Birch} and {Swinnerton-Dyer} and a geometric analog.
\newblock In \emph{S\'eminaire Bourbaki, Vol.~9}, pages Exp.~No.~306, 415--440. Soc. Math. France, 1966.

\bibitem[{The LMFDB Collaboration}(2024)]{LMFDB}
{The LMFDB Collaboration}.
\newblock The {L}-functions and modular forms database.
\newblock \url{https://www.lmfdb.org}, 2024.
\newblock Online database.

\bibitem[Tsiokos(2026{\natexlab{a}})]{TsiokosBSDClosure2026}
Ioannis Tsiokos.
\newblock {A Conditional Proof of the Strong Birch--Swinnerton-Dyer Conjecture}, 2026{\natexlab{a}}.
\newblock Preprint. \doi{10.5281/zenodo.20713968}.

\bibitem[Tsiokos(2026{\natexlab{b}})]{TsiokosFoundationsII2026}
Ioannis Tsiokos.
\newblock {Six Birds Foundations II: Admissibility Meta-Theory and the Exact Six Program}, 2026{\natexlab{b}}.

\bibitem[Tsiokos(2026{\natexlab{c}})]{TsiokosFoundationsIII2026}
Ioannis Tsiokos.
\newblock {Six Birds Foundations III: A Finite Audited Interaction Calculus for SBT}, 2026{\natexlab{c}}.

\bibitem[Tsiokos(2026{\natexlab{d}})]{TsiokosNeedles2026}
Ioannis Tsiokos.
\newblock {Six Birds: The Needles Framework Apparatus (adequacy, Schur-complement diagnostics)}, 2026{\natexlab{d}}.

\bibitem[Wiles(1995)]{Wiles1995}
Andrew Wiles.
\newblock Modular elliptic curves and {Fermat's} last theorem.
\newblock \emph{Ann. of Math. (2)}, 141\penalty0 (3):\penalty0 443--551, 1995.
\newblock \doi{10.2307/2118559}.

\bibitem[Zhang(2005)]{Zhang2005schur}
Fuzhen Zhang, editor.
\newblock \emph{The {Schur} Complement and Its Applications}, volume~4 of \emph{Numerical Methods and Algorithms}.
\newblock Springer, 2005.
\newblock \doi{10.1007/b105056}.

\end{thebibliography}

\endgroup
\end{document}
